\section{可換復元の定式化と主結果}\label{sec:intro}

滑らかなcubicのKuznetsov成分を
\[
 D^b(X)=\langle\Ku(X),\OO_X,\OO_X(1),\ldots,\OO_X(n-2)\rangle
\]
の規約で取る。rotationは
\[
 \rho_X=\mathbf L_{\OO_X}\circ(-\otimes\OO_X(1))
\]
である。すべての圏はdg enhancementと冪等完備化を保って扱う。
可換cubicについては、\((\Ku(X),\rho_X)\) が \(X\) の射影同型類を決定する
\cite[Theorem 1.7]{LZ}。このTorelli定理は可換モデルの一意性を与えるが、
任意の非可換な圏の組がその像に属することを述べない。
幾何側のenhanced同値は、包含・射影と合成してFourier--Mukai核で表されるため、
同定理のFourier--Mukai条件を満たす。
本稿は、この存在問題の形式近傍、飽和な代数的族のétale近傍、
および明示的な楕円ブロックの族を扱う。

\subsection{形式的な存在と一意性}

\begin{theorem}[次数付き変形の形式的可換モデル]\label{thm:main}
 \(X=V(F)\subset\PP^{n+1}_{\CC}\) を滑らかなcubic、\(n\ge3\) とする。
 \(R\) を剰余体が \(\CC\) の有限次元可換局所Artin代数とし、
 \(B_R=\bigoplus_{d\ge0}(B_R)_d\) を中心的な \(R\)-代数で、
 \((B_R)_0=R\)、各次数部分が有限自由、かつ指定された次数付き同型
 \[
 B_R\otimes_R\CC\simeq\CC[x_1,\ldots,x_{n+2}]/(F)
 \]
 を持つものとする。次を定義する。
 \[
 \cP_R=D(\QGr B_R)^c,\qquad
 \cK_R=\langle\OO_R,\ldots,\OO_R(n-2)\rangle^\perp\subset\cP_R,\qquad
 \rho_R=\mathbf L_{\OO_R}\circ(1).
 \]
 このとき \((\cK_R,\rho_R)\) はsmooth properな \(R\)-linear dg圏と自己同値の組である。
 あるmarkedな滑らかなcubic変形
 \(X_R\subset\PP^{n+1}_R\) と、markingを保つ同値
 \[
 (\cK_R,\rho_R)\simeq(\Ku(X_R),\rho_{X_R})
 \]
 が存在する。\(X_R\) のmarked射影同型類は一意である。
 これらはArtin底の変更と両立する。
\end{theorem}

ここで \(\QGr\) は全次数付き加群をtorsion加群で割ったGrothendieck圏であり、
上付き \(c\) はcompact対象を表す。Artin底上では
\(\cP_R\) を \(D^b(\qgr B_R)\) 全体に置き換えない。
例えば可換な定数族でも、底 \(R\) が非正則なら両者は異なる。
後述する体上の有限パラメータでは、残余圏を通常のgraded matrix factorization圏、
すなわちOrlov同値で \(D^b(\qgr B)\) に埋め込まれる成分として扱う。
特異なファイバーで、この成分と \(D(\QGr B)^c\) の直交成分を混同しない。

\begin{corollary}[形式族]\label{cor:formal}
 \(B_\hbar\) が次数ごとに \(\CC[[\hbar]]\)-flatな、定理の斉次座標環の形式的変形なら、
 \[
 F_\hbar\in\CC[[\hbar]][x_1,\ldots,x_{n+2}]_3,\qquad F_\hbar\bmod\hbar=F
 \]
 が存在し、各 \(\CC[\hbar]/(\hbar^{r+1})\) 上で残余圏の組を
 \(V(F_\hbar)\) の組と同一視できる。これらの同値は互換的に選べる。
 形式的cubicのmarked同型類は一意である。
\end{corollary}

これは形式的な主張である。係数の収束や、指定された代数的パラメータ空間への
降下、任意の非零パラメータでの可換モデルの存在を含まない。
また \(F_\hbar=F\) を意味しない。

\subsection{代数的族における局所的存在}

定理\ref{thm:etale}は、有限型複素底 \(S\) 上のsmooth properな代数的dg圏の族
\(\cK_S\) と、可逆perfect bimoduleで表された自己同値 \(\rho_S\) を入力とする。
ある点 \(s\) で \((\cK_s,\rho_s)\simeq(\Ku(X),\rho_X)\) なら、
\(s\) のétale近傍 \(T\) 上に滑らかなcubic \(Y/T\) と
\[
 (\cK_T,\rho_T)\simeq(\Ku(Y/T),\rho_{Y/T})
\]
が存在する。従って、可換cubicで実現される複素点はZariski開部分をなす。
これはgraded flatな代数の族から、相対smooth properなdg模型を作る課題を
省略するものではない。

また \(n\ne4\) では \(\HH^2(\Ku(X))=J(F)_3\)、\(\HH^1=0\) であり、
rotationを忘れた形式的復元も成立する（定理\ref{thm:unpolarized-formal}）。
同じ条件下では、代数的dg圏の族についてもrotationを最初から指定せずに
étale局所モデルを作れる（系\ref{cor:etale-unpolarized}）。
これらは局所的な実現の結果であり、rotationなしの大域的な一意性とは分けて扱う。

\subsection{有限パラメータを含む明示的族}

\[
 S_h=\CC\langle x,y,z\rangle/
 ([x,y]-hz^2,\ [y,z]-hx^2,\ [z,x]-hy^2)
\]
と置き、\(T=\sym(xyz)\) を六つの順列の平均、
\(G_h=x^3+y^3+z^3-2h^2T\) とする。
中心変数 \(u_1,\ldots,u_m\) と斉次三次式 \(P(u)\) に対して
\[
 B_{h,P}=S_h[u_1,\ldots,u_m]/(G_h+P(u))
\]
を考える。\(P=\sum_i u_i^3\) の場合を \(B_{h,m}\) と略記する。

\begin{theorem}[一ブロックと任意の中心三次式の可換モデル]\label{thm:global}
 \(m\ge2\)、\(n=m+1\)、\(h\in\CC\)、\(h^6\ne-27\) とする。
 \(P\in\CC[u_1,\ldots,u_m]_3\) のアフィン臨界点が原点だけであると仮定する。
 \(B_{h,P}\) の残余dg圏とrotationは、滑らかなcubic
 \[
 Y_{h,P}=V\!\left(x^3+y^3+z^3+h^2xyz+P(u)\right)
 \subset\PP^{m+2}
 \]
 の \((\Ku(Y_{h,P}),\rho_{Y_{h,P}})\) と同値である。
 この \(Y_{h,P}\) は、同じrotation付き残余圏を持つ滑らかな可換cubicとして
 射影同型を除いて一意である。
\end{theorem}

この定理はsixfoldでは \(m=5\)、sevenfoldでは \(m=6\) を含む。
中心変数どうしの任意の三次混合項を許すが、一つのSklyaninブロックと中心部分が
分離されているという仮定を保持する。\(P=\sum u_i^3\) はその特殊例である。
複数の非可換ブロックについても、対応する可換cubicの残余圏からの完全忠実埋込みを得る
（命題\ref{prop:multiblock-embedding}）。次の部分的な有限位数条件の下では、その像が全体を生成する。

\begin{theorem}[一つの任意位数ブロックを許す多ブロック可換モデル]\label{thm:multiblock-model}
 \(r\ge2\)、\(m\ge0\)、\(n=3r+m-2\) とする。
 各 \(j\) について \(h_j^6\ne-27\) とし、楕円曲線
 \[
 E_{h_j}=V(x_j^3+y_j^3+z_j^3+h_j^2x_jy_jz_j),\qquad
 o_j=[1:-1:0],\quad p_j=[1:-1:-h_j]
 \]
 上の点 \(p_j\) は、高々一つを除いて有限位数であると仮定する。
 すなわち因子を並べ替えた後、\(j\ge2\) にのみ有限位数を仮定し、\(p_1\) の位数には条件を課さない。
 \(m>0\) の場合、\(P\in\CC[u_1,\ldots,u_m]_3\) は原点だけに臨界点を持つとし、
 \(m=0\) の場合は中心部分を省く。
 異なるSklyaninブロックの元を可換とした代数
 \[
 B_{\mathbf h,P}=
 (S_{h_1}\otimes\cdots\otimes S_{h_r})[u_1,\ldots,u_m]
 \big/\left(\sum_jG_{h_j}+P\right)
 \]
 のrotation付き残余dg圏は、滑らかなcubic
 \[
 Y_{\mathbf h,P}=V\!\left(
 \sum_{j=1}^r(x_j^3+y_j^3+z_j^3+h_j^2x_jy_jz_j)+P(u)
 \right)\subset\PP^{3r+m-1}
 \]
 の \((\Ku(Y_{\mathbf h,P}),\rho_{Y_{\mathbf h,P}})\) と同値である。
 この可換モデルは射影同型を除いて一意である。
\end{theorem}

ここで「有限パラメータ」は形式変数でなく実際の複素数であることを意味し、
楕円曲線上の点が「有限位数」であることとは異なる。
定理は一つの無限位数ブロックを許す。
具体的に \(h_1=2\) は無限位数（補題\ref{lem:order-infinite}）、
\(h_j^6=5\ (j\ge2)\) は位数五（補題\ref{lem:order-five}）となる。
従って二ブロックのsixfold、三ブロックのsevenfoldにも、真に無限位数の因子を含む例が得られる。
二つ以上の無限位数ブロックを任意に許す主張は、本定理には含まれない。
証明は第\ref{sec:finite-center}節に置く。

\begin{proposition}[例外パラメータのhomological singularity]\label{prop:bad-intro}
 \(m\ge0\) とする。Fermat中心部分の族 \(B_{h,m}\) で \(h^6=-27\) とすると、\(S_h[u]\) はAS-regularで
 \(G_h+\sum u_i^3\) は中心正則元であるが、そのgraded matrix factorization圏は
 properではなく、さらに
 \[
 \operatorname{hdim}(\qgr B_{h,m})=\infty
 \]
 である。ここでhomological dimensionは、非零のExtが存在する次数の上限である。
 従って、滑らかな射影cubicのKuznetsov成分とは同値にならない。
\end{proposition}

これはAS-regularな周囲、三次の中心正則元、期待されるHilbert series、
および \(\rho^3\simeq[2]\) だけでは、大域的な可換モデルの存在を保証しないことを示す。
本稿の可換点からのPoisson量子化族に限定した、良い有限パラメータ全体での存在問題は別に残る。
一般のAS-regular中心三次商については、smooth properかつCalabi--Yauな例にも
定理\ref{thm:zhang-cy-fourfold-counterexample}の反例がある。

\subsection{先行研究と本稿の位置}

Orlov \cite{Orlov} とCassidy--Conner--Kirkman--Moore \cite{CCKM} は、
中心正則元による商、特異圏、matrix factorization、残余成分の比較を与える。
体上のwindowのdg模型にはCăldăraru--Tu \cite[Theorem 5.8]{CT} を用いる。
本稿の形式的な構成では、これらの体上の同値をArtin底上へ無条件に延長せず、
De Deken--Lowen \cite{DDL} のabelian deformationとcompact generationを用いる。
smoothnessとpropernessの持ち上げはnilpotent descent
\cite{Mathew}による。

残余圏のHochschild cohomologyはBallard--Favero--Katzarkov \cite{BFK}、
幾何学的Kodaira--Spencer写像はHuybrechts--Rennemo \cite{HR} と、
三次元の場合にはJacovskis--Lin--Liu--Zhang \cite{JLLZ} を用いる。
四次元の余分な二方向に対するrotationの作用は
Huybrechts \cite[Proposition 3.12]{Huy} から確認する。
actual dg変形をHochschild形式変形と同一視する際のcurvature問題
\cite{KL}は、Lurie \cite{Lurie} とBlanc--Katzarkov--Pandit
\cite{BKP} の比較写像の単射性で処理する。

明示的族の楕円ブロックはEtingof--Ginzburg \cite{EG} と
Sklyaninの点曲線の公式 \cite{Walton} に依拠する。
本稿では次数シフトの共役、中心変数を加える際の生成性、
有限群による通常次数への降下を証明する。
可換な係数環についてのThom--Sebastiani定理を、
非可換な係数環に対して仮定を確認せず引用することはしない。

有限中心代数の外部積生成については、Elagin--Lunts--Schn\"urer
\cite[Theorem 4.18]{ELS} が、perfect field上のnoetherian J-2 schemesに載る
coherentな代数に対し、ungradedな有界導来圏の生成定理を証明している。
本稿では、\(\ZZ^r\) の任意の粗化次数を保ち、さらに一方の係数代数には
有限中心性を仮定しない形を証明する（定理\ref{thm:one-sided-boxes}）。
\cite[Theorem 4.18]{ELS} は双方の因子をcoherentな構造代数とするため、
任意のnoetherian係数代数を許すここでの主張は、その定理の直接の適用ではない。
新たな入力はBeilinsonの対角分解 \cite{Beilinson} を非可換係数付きで用いる
多項式拡大の補題である。中心のNoether正規化と分離冪等元によるretractを組み合わせ、
斉次中心局所化の核を次元について帰納する。
対象の細分次数へのliftや、係数代数の有限global dimensionは仮定しない。
有限位数THCRの有限中心性 \cite[Corollary 2.3]{SmithTate} をこの片側のみに適用し、
一つの任意位数ブロックを含む多ブロックのMF比較を得る。
He--Ma--Ye \cite{HMY} のtensor productに対する一般化Kn\"orrer周期性は、
中心元の次数が二である非可換quadricとClifford変形を扱う。
本稿の三次中心元への直接の適用ではない。
周囲のnoetherian性については、Chan \cite{ChanMulti} の多重Serre同値と
面・帯による議論を任意個の楕円因子へ展開し、有限位数の仮定なしで証明する。
Keeler \cite{KeelerMulti} のTheorem 2.7は任意個の可換な可逆bimoduleの
NC-amplenessを判定するが、同論文のTheorem 3.3とCorollary 3.5のnoetherian性は
二因子についての主張である。本稿では任意個の積に必要な帰納法を記す。
基礎となる対角分解と有限中心のungraded生成定理は既知である。
一方有限中心の次数付き定式化および本適用の先行研究上の新規性は確認継続中である。

étale局所復元では、Toën--Vaquié \cite{TV} のperfect objectのmoduliと
smoothnessの無限小判定 \cite{Stacks} を用いる。
新たに示すのは、同値bimoduleとrotation intertwinerを同時にパラメータ化したスタックが、
形式的復元定理によって底上smoothになることである。
Pirozhkov \cite{Pirozhkov} およびLin--Zhang \cite{LZ} のTorelliは
二つの可換モデルの一意性に関する結果であり、この局所的な存在の証明とは役割が異なる。
Huybrechts \cite[Section 9]{HuyUpdate} の四次元での一般的な非可換変形の問題提起と、
本稿のrotation付きの変形問題との接写像の詳細な比較は研究ノートに残す。

\paragraph{主張の状態.}
本稿の定理には以下で証明を与える。これらを一つの一般次元の復元定理として
整理したことの先行研究上の新規性は未確認であり、
既知の基礎定理自体を新しい成果として扱わない。
本稿で構成する可換点からのPoisson量子化族について、
良い有限パラメータ全体での存在は未解決として残す。

\section{次数付き変形から相対残余圏を作る}\label{sec:relative}

\subsection{nilpotentな底変更によるcompactnessの検出}

\begin{lemma}\label{lem:descent}
 \(R\) を上の局所Artin代数、\(k=R/\fm\) とする。
 \(A\) を \(R\)-linear dg代数、\(M\) を \(A\)-加群とし、
 \(A_0=A\otL_R k\)、\(M_0=M\otL_R k\) と書く。
 \(M_0\) が \(A_0\)-compactなら \(M\) は \(A\)-compactである。
 また導来底変更 \(-\otL_R k\) は保守的である。
\end{lemma}
\begin{proof}
 \(\fm\) の冪による有限濾過から
 \(R\in\thick_{D(R)}(k)\) である。
 \(A\)-加群の任意の族 \((N_\alpha)\) に対し、
 \[
 \bigoplus_\alpha\RHom_A(M,P\otL_RN_\alpha)
 \longrightarrow
 \RHom_A\!\left(M,P\otL_R\bigoplus_\alpha N_\alpha\right)
 \]
 が同値となる \(P\in D(R)\) の全体はthickである。
 \(P=k\) では導来随伴と \(M_0\) のcompactnessによって同値となるので、
 \(P=R\) でも同値となる。保守性も、
 \(P\otL_RM=0\) となる \(P\) のthick部分圏を使って同様に従う。
 これは \cite[Propositions 3.28, 3.35; Remark 3.29]{Mathew}
 の特別な場合である。
\end{proof}

\begin{remark}
 この補題は \(k\) が \(R\)-perfectであるとは仮定していない。
 また普通の商 \(A/\fm A\) だけを調べる命題ではなく、
 必ず導来ファイバー \(A\otL_Rk\) を使う。
\end{remark}

\subsection{abelian deformationと明示的な生成子}

\begin{proposition}\label{prop:qgr}
 定理\ref{thm:main}の \(B_R\) について、
 \(\QGr B_R\) は \(\operatorname{Qcoh}(X)\) のflat abelian deformationである。
 \[
 E_R=\bigoplus_{i=0}^{n+1}\OO_R(i)
 \]
 は \(D(\QGr B_R)\) のcompact generatorであり、
 \[
 \Lambda_R=\REnd(E_R),\qquad
 \Lambda_R\otL_Rk\simeq\REnd_X(E_0)
 \tag{\(\ast\)}\label{eq:derived-basechange}
 \]
 がdg代数として成り立つ。
\end{proposition}
\begin{proof}
 まず \(B_R\) は左右noetherianである。
 \(\fm^aB_R/\fm^{a+1}B_R\) は、
 flatnessにより \(B_0\) の有限個の直和である。
 有限濾過と \(B_0\) のnoetherian性を使えばよい。
 次数付き生成元を持ち上げる同じ議論から、
 \(B_R\) は有限生成で、付随する \(\ZZ\)-代数のtailsも通常のtorsion商で記述できる。

 \(X\) の射影正規性により \(B_0\) の \(\ZZ\)-代数は
 \((\OO_X(i))_{i\in\ZZ}\) のdirected Hom-categoryである。
 この列は \(\ZZ\)-generatingであり、\(i\le j\) に対して
 \[
 \Ext_X^{1,2}(\OO_X(i),\OO_X(j))
 =H^{1,2}(X,\OO_X(j-i))=0.
 \]
 hypersurfaceの完全列と射影空間のcohomologyからこの消滅が従う。
 よって \cite[Theorem 4.5]{DDL} を適用できる。
 \(\OO_R(i)\) は明示的なflat liftである。
 実際 \(B_R(i)\) との有限 \(R\)-加群のテンソル積は完全で、
 \(\operatorname{Gr}B_R\to\QGr B_R\) は完全である。
 これらの対象はabelian圏も生成するので、必要なflat resolutionを取れる。

 閉埋め込み \(j:X\hookrightarrow\PP^{n+1}\) を考える。
 \(E_0=j^*(\bigoplus_{i=0}^{n+1}\OO_{\PP^{n+1}}(i))\) はcompactである。
 \(\RHom_X(E_0,M)=0\) なら、随伴とBeilinsonの生成定理
 \cite{Beilinson}から \(Rj_*M=0\)、従って \(M=0\) となる。
 したがって \(E_0\) はcompact generatorである。

 \(E_R\) は有界flat liftなので、
 \cite[Propositions 5.8, 5.9]{DDL} によりcompact generatorとなる。
 同論文Proposition 5.5と導来随伴は
 \[
 k\otL_R\RHom(E_R,E_R)
 \simeq\RHom_X(k\otL_RE_R,k\otL_RE_R)
 \]
 を与える。同論文Theorem 5.13の証明でも同じ底変更が使われている。
 dg enhancement上の導来還元関手からこの比較を作るため、
 比較は合成と単位を保つ。必要ならhomotopically \(R\)-flatなdg模型を選ぶ。
 これで \eqref{eq:derived-basechange} は複体だけでなくdg代数として成立する。
\end{proof}

\begin{corollary}\label{cor:smooth-proper}
 \(\cP_R=D(\QGr B_R)^c\) はsmooth properであり、
 \(\cP_R\otL_Rk\simeq\Perf(X)\) である。
\end{corollary}
\begin{proof}
 導来Morita理論 \cite{Keller} によって
 \(\cP_R\simeq\Perf(\Lambda_R)\) である。
 \(X\) のpropernessから \(\Lambda_0\) は \(k\)-perfectなので、
 補題\ref{lem:descent}で \(\Lambda_R\) は \(R\)-perfectとなる。
 次に
 \[
 A=\Lambda_R^{\mathrm{op}}\otL_R\Lambda_R,\qquad M=\Lambda_R
 \]
 とする。この導来ファイバーは \(\Lambda_0\) のenveloping algebraと対角加群である。
 \(X\) のsmoothnessにより対角加群はその上でperfectであり、
 同じ補題から \(\Lambda_R\) の対角もperfectとなる。
\end{proof}

\subsection{例外列、mutation、rotation}

\begin{proposition}\label{prop:relative-pair}
 \(\OO_R,\ldots,\OO_R(n-2)\) は \(\cP_R\) 内のrelative exceptional collectionである。
 その右直交 \(\cK_R\) はsmooth properで、導来特殊ファイバーは \(\Ku(X)\) である。
 \(\rho_R=\mathbf L_{\OO_R}\circ(1)\) は \(\cK_R\) の自己同値となる。
 これらの構成はArtin底変更と両立する。
\end{proposition}
\begin{proof}
 Hom複体の導来底変更と補題\ref{lem:descent}から
 \[
 \RHom(\OO_R(i),\OO_R(i))\simeq R,\qquad
 \RHom(\OO_R(j),\OO_R(i))=0\quad(j>i,\ 0\le i,j\le n-2)
 \]
 である。自己Homについては単位写像のconeを還元する。
 propernessにより評価・余評価はperfect対象を保つため、各例外成分はadmissibleである。
 評価のconeによるmutationを順に取って
 \[
 p_R=\mathbf L_{\OO_R}\cdots\mathbf L_{\OO_R(n-2)}
 \]
 と置くと、\(G_R=p_R(E_R)\) は \(\cK_R\) を生成する。
 mutationの導来底変更は \cite[Lemma 5.14]{DDL} による。
 \(G_R\) の導来特殊ファイバーは \(\Ku(X)\) の生成子なので、
 その自己準同型dg代数へ補題\ref{lem:descent}を適用すればsmooth proper性も従う。

 \(T=(1)\) は次数付き射影圏の自己同値で、\(T\OO_R(i)=\OO_R(i+1)\) である。
 \(M\in\cK_R\)、\(1\le i\le n-2\) に対して
 \[
 \RHom(\OO_R(i),TM)=\RHom(\OO_R(i-1),M)=0.
 \]
 \(\OO_R\) によるmutationはこの消滅を保ち、\(i=0\) の直交性も与える。
 よって \(\rho_R\) は \(\cK_R\) を保つ。
 その特殊ファイバーは通常の \(\rho_X\) である。
 Homの比較写像は還元後に同値なので、nilpotentな保守性から
 \(\rho_R\) は完全忠実である。

 \(\rho_R(G_R)\) の特殊ファイバーは生成子である。
 これに直交するInd対象を導来還元すると零になり、
 同じ保守性で元の対象も零になる。従って \(\rho_R(G_R)\) もcompact generatorである。
 冪等完備な圏からの完全忠実完全関手の像はthickなので、\(\rho_R\) は本質的全射でもある。

 底変更については、次数付き導来テンソル関手による生成子の自己Homの自然な比較を使う。
 剰余体上でこの比較は恒等写像となるため、変更後のArtin底上で同値となる。
 評価写像とそのconeも底変更と両立するので、残余圏とrotationまで両立する。
\end{proof}

\section{rotation付きの形式的復元}\label{sec:formal}

\subsection{正確な変形問題}

\begin{definition}
 \((\Ku(X),\rho_X)\) のmarked pair deformationとは、
 smooth properな \(R\)-linear dg圏 \(\cK_R\)、invertible bimoduleで表される
 自己同値 \(\rho_R\)、および特殊ファイバーでのenhanced pairの指定された同定からなる。
 同値にはこのmarkingとの両立を要求する。
 \(\rho_R^3\simeq[2]\) の指定されたcoherenceは定義に含めない。
\end{definition}

\begin{theorem}[pairの形式的復元]\label{thm:pair}
 \(X\) を滑らかなcubic \(n\)-fold、\(n\ge3\) とする。
 marked cubic deformationからmarked pair deformationへの写像は、
 各古典的局所Artin \(\CC\)-代数上で同型類の全単射を与える。
 互換的な形式族についても同型類の一意性と互換的な実現が成り立つ。
\end{theorem}

この主張はfull higher deformation groupoidの同値ではない。
pairには自然変換のscalar automorphismがあり、実現する圏同値も一意とは限らない。

\subsection{actual変形とHochschild形式変形}

\(\cC=\Ind(\Ku(X))\) とする。
actual compactly generated deformationの空間と、その形式moduli problemの間には
\[
 \theta_{\cC}:\operatorname{CatDef}_{\cC}^{c}
 \longrightarrow\operatorname{CatDef}_{\cC}^{\wedge}
\]
がある。右辺の接複体は \(CC^\bullet(\Ku(X))[2]\) である。
actual Morita変形が任意のHochschild classを実現するとは一般にはいえない
\cite{KL}。本稿ではその全射性を使わない。

\begin{lemma}\label{lem:theta}
 上の \(\theta_{\cC}\) は \((-1)\)-truncatedである。
 従って、右辺で同値となった二つのactual変形の同値はactual側へ戻せる。
 この性質はrotationによる \(\ZZ\)-homotopy fixed pointを取った後にも成り立つ。
\end{lemma}
\begin{proof}
 \(\Ku(X)\) はproperでcompact generatorを持つためtameである。
 \cite[Proposition 5.3.21, Theorem 5.1.9]{Lurie} を適用する。
 \cite[Definition 3.3, Section 3.3]{BKP} にも、
 \(\pi_0\) で単射、\(\pi_i\), \(i\ge1\), で同型という形でこの帰結が明記されている。
 比較写像は圏同値に自然であり、markingを移すrotationの作用と両立する。
 空間のmonomorphismはhomotopy limitで保たれる。
\end{proof}

一つのenhanced autoequivalenceを持ち上げる問題は、
markingに作用する \(\ZZ\) のhomotopy fixed pointである。
\(\ZZ\) は一つの可逆射で自由生成されるので、追加の有限群関係を課していない。
\[
 P=\left(\operatorname{CatDef}_{\cC}^{\wedge}\right)^{h\ZZ}
\]
と置くと、これはformal moduli problemであり、
\[
 T_P=R\Gamma(\ZZ,CC^\bullet(\Ku(X))[2]).
\]
\(\ZZ\) の自明加群の長さ一の自由分解から、次の完全列を得る。
\[
 0\longrightarrow
 \operatorname{coker}(1-\rho:\HH^1\to\HH^1)
 \longrightarrow H^0(T_P)
 \longrightarrow
 \ker(1-\rho:\HH^2\to\HH^2)
 \longrightarrow0.                                      \tag{1}\label{eq:tangent}
\]

\subsection{Hochschild計算とKodaira--Spencer写像}

\(N=n+2\) および
\[
 J(F)=\CC[x_1,\ldots,x_N]/(\partial_1F,\ldots,\partial_NF)
\]
を用いる。

\begin{lemma}\label{lem:HH}
 \(n\ge3\) のとき
 \[
 \HH^{<0}(\Ku(X))=0,\quad \HH^0(\Ku(X))=\CC,\quad
 \HH^1(\Ku(X))=0,
 \]
 \[
 \HH^2(\Ku(X))=
 \begin{cases}
 J(F)_3& n\ne4,\\
 J(F)_3\oplus\CC^2& n=4.
 \end{cases}
 \]
 幾何学的Kodaira--Spencer写像は同型
 \[
 H^1(X,T_X)\xrightarrow{\sim}
 \HH^2(\Ku(X))^{\rho_X}=J(F)_3
 \]
 を与える。その次元は \(\binom{n+2}{3}\) である。
\end{lemma}
\begin{proof}
 \cite[Theorem 5.39]{BFK} のKoszul cohomologyによる公式を
 三次のequivariant matrix factorizationへ適用する。
 恒等sectorは偶数次数 \(2r\) に \(J(F)_{3r}\) を与える。
 非自明な \(g\in\mu_3\) では固定点集合は原点だけであり、
 conormalのrankは \(N\)、determinant characterは \(-N\) である。
 cohomological degree \(q\) で残るweightは
 \[
 N+3(q-N)/2=(3q-N)/2.
 \]
 従って二つの非自明sectorは、\(3\mid N\) のときに限り、
 \(q=N/3\) に各一次元を与える。これで列記した群が得られる。

 BFKの恒等sectorのKoszul模型では、\(G\) の乗法が \([G]\) を表す。
 \(F+\varepsilon G\) の実際の変形類との一致は次の比較による。
 \(n\ge4\) では、periodicityから作るkernel morphismと
 Atiyah--Kodaira--Spencer morphismの一致を
 \cite[Lemma 2.19]{HR} とその前後の構成で確認できる。
 ここで三次に必要な次数条件 \(3\le(n+2)/2\) は満たされる。
 \(n=3\) では \cite[Theorem 4.4(2)]{JLLZ} が必要な接写像の同型を与える。
 normal sequenceとEuler sequenceは \(H^1(T_X)=J(F)_3\) を与える。

 幾何学的変形はrotationを持つので、この部分空間は不変である。
 \(n\ne4\) ではそれが \(\HH^2\) 全体である。
 \(n=4\) では \cite[Proposition 3.12]{Huy} により、
 rotationはMukai格子の \(A_2^\perp\) を固定し、\(A_2\) の三つの根を巡回置換する。
 invariantな2-Calabi--Yau volumeによる
 \(\HH^2\simeq\HH_0\) の下で、余分な二方向の固有値は
 \(\zeta,\zeta^2\) となり、不変部分はない。
 最後に偏微分は二次の正則列であるから
 \(H_{J(F)}(t)=(1+t)^N\) であり、次数三の次元は \(\binom N3\) である。
\end{proof}

\subsection{持ち上げの差と自己同型}

\begin{lemma}[Artin帰納法]\label{lem:induction}
 \(D\) をsmoothかつinfinitesimally homogeneousなmarked変形関手で、
 古典的Artin底上ではset-valuedとし、
 \(P\) をformal moduli problemとする。
 自然な写像 \(f:D\to P\) が接空間で同型を誘導すると仮定する。
 さらに、各small extension \(R'\to R\) と \(x'\in D(R')\) に対し、
 \(f(x'\vert_R)\) の任意の自己同型類が \(f(x')\) 上へ持ち上がると仮定する。
 このとき \(D(R)\to\pi_0P(R)\) は各古典的Artin底上で全単射である。
\end{lemma}
\begin{proof}
 Artin長について帰納法を行う。核 \(I\) のsmall extension上に
 \(y'\in P(R')\) があり、その還元を \(f(x)\) と同一視できたとする。
 \(D\) のsmoothnessにより \(x\) を \(x'\) に持ち上げる。
 固定した還元の同定を持つ \(f(x')\) と \(y'\) は、
 \(P(R')\to P(R)\) の同じhomotopy fibreに属する。
 このfibreが空でなければ、その連結成分は
 \(H^0(T_P)\otimes I\) のtorsorである。
 接写像の全射性によって \(x'\) を修正すれば差を消せる。

 単射性では、\(f(x')\simeq f(z')\) の還元を考える。
 帰納法により幾何学的還元は一致するが、
 圏同値の還元はその像の自己同型となり得る。
 仮定によってこの自己同型を上へ持ち上げ、その逆を合成する。
 すると還元を恒等に固定したfibreで二つのliftを比較できる。
 差の接ベクトルの像は零であり、接写像の単射性で元の差も零となる。
\end{proof}

\begin{lemma}\label{lem:automorphisms}
 cubicから得たpairのmarked deformationにおいて、
 small extensionに沿う上の自己同型持ち上げ条件が成り立つ。
\end{lemma}
\begin{proof}
 smooth properな圏ではHochschild複体はperfectで、その導来底変更が成り立つ。
 補題\ref{lem:HH}と局所Artin環上のminimal perfect complexを使うと
 \[
 \HH^0(\cK_R/R)=R,\qquad \HH^1(\cK_R/R)=0
 \]
 である。実際特殊ファイバーの負次数と次数一が零なので、
 minimal complexにはそれらの次数の項がなく、次数零は単位で生成される。

 恒等関手に還元されるinvertible bimoduleは相対対角と同型である。
 perfect bimoduleのliftの差はsmall extensionの核 \(I\) に対して
 \(\Ext^1(\Delta,\Delta)\otimes I=\HH^1(\Ku(X))\otimes I=0\)
 で測られ、対角自身がliftを与えるためである。
 この同型は各段階で指定した特殊ファイバーのmarkingを保つように選ぶ。
 従ってpairの自己同型の基礎となる圏同値は恒等と同型であり、
 残るのは \(\rho_R\Rightarrow\rho_R\) のscalar unitである。
 markingを固定すればこのscalarは \(1+\fm_R\) に入り、必ず \(R'\) へ持ち上がる。
 scalarな二次自己同型も持ち上がる。
 formal moduli problem側の自己同型をactual bimoduleで扱えることは、
 補題\ref{lem:theta}の \(\pi_1\) の同型による。
\end{proof}

\begin{proof}[定理\ref{thm:pair}の証明]
 cubicのmarked変形関手はsmoothである。方程式の係数を持ち上げればよく、
 \(H^1(\OO_X)=H^2(\OO_X)=0\) によって \(\OO_X(1)\) は一意に延長する。
 infinitesimalな射影自己同型もない。実際、二次の偏微分の正則列には
 一次係数の非自明なsyzygyがなく、Euler scalarを除けば \(H^0(T_X)=0\) である。
 これはmarkedな幾何学的変形をset-valuedとして扱うことを正当化する。

 式\eqref{eq:tangent}と補題\ref{lem:HH}により接写像は同型である。
 補題\ref{lem:automorphisms}と補題\ref{lem:induction}を適用し、
 最後に補題\ref{lem:theta}でactual pairの同値へ戻す。
 帰納法では各段階で指定した還元の同値を保てるので、
 \(\CC[\hbar]/(\hbar^{r+1})\) に沿った実現も互換的に選べる。
\end{proof}

\begin{proof}[定理\ref{thm:main}と系\ref{cor:formal}の証明]
 命題\ref{prop:relative-pair}の組へ定理\ref{thm:pair}を適用する。
 互換的な幾何学的liftの方程式の係数の逆極限が \(F_\hbar\) を与える。
 この最後の操作は形式族を作るものであり、非零複素数で評価する操作ではない。
\end{proof}

\begin{corollary}\label{cor:period}
 上で構成したArtin pairについて、指定を伴わない同型
 \(\rho_R^3\simeq[2]\) が存在する。
\end{corollary}
\begin{proof}
 \(\rho_R^3[-2]\) は恒等に還元されるinvertible bimoduleである。
 補題\ref{lem:automorphisms}で使った \(\HH^1\)-rigidityを適用する。
 または可換モデルからこの同型を移してもよい。
 特定のmatrix factorization表示が与えるperiodicity同型とのcoherenceの一致は、
 この系の主張に含めない。
\end{proof}

\begin{remark}
 sevenfoldでは \(\HH^3(\Ku(X))\simeq\CC^2\) があるが、
 上の帰納法は \(\HH^3=0\) を要求しない。
 比較している二つのliftは既に存在し、その差を接空間で修正する。
 higher operationsを無視したformalityの仮定も使っていない。
\end{remark}

\subsection{四次元以外ではrotationを忘れた形式復元}\label{sec:unpolarized-formal}

四次元以外では、形式近傍に限れば、復元の入力からrotationを省くことができる。
ここでいうmarked dg変形は、特殊ファイバーを \(\Ku(X)\) と同定した
smooth properな \(R\)-linear dg圏であり、同値にもこのmarkingとの両立を要求する。

\begin{theorem}[rotationを忘れた形式復元]\label{thm:unpolarized-formal}
 \(X\subset\PP^{n+1}_{\CC}\) を滑らかなcubicとし、\(n\ge3\)、\(n\ne4\) とする。
 任意の古典的局所Artin \(\CC\)-代数 \(R\) について、
 marked cubic変形から \(\Ku(X)\) のmarked smooth proper dg変形への写像は、
 同型類の全単射を与える。従って、任意のそのような \(\cK_R\) は
 一意なmarked cubic変形 \(X_R\) の \(\Ku(X_R)\) と同値である。

 また、固定した \(\cK_R\) 上で、\(\rho_X\) を持ち上げるinvertible bimodule
 \(\rho_R\) が存在し、そのmarked同型類は一意である。
 従ってrotationを忘れる写像
 \[
  \{\text{marked smooth proper pair変形}\}/\simeq
  \ \longrightarrow\
  \{\text{marked smooth proper dg変形}\}/\simeq
 \]
 も全単射である。これらの主張はArtin底変更と互換的な形式族についても成り立つ。
\end{theorem}
\begin{proof}
 \(Q=\operatorname{CatDef}_{\Ind\Ku(X)}^{\wedge}\) と置く。
 補題\ref{lem:HH}の幾何学的Kodaira--Spencer写像は、\(n\ne4\) では
 \[
  H^1(X,T_X)\xrightarrow{\sim}
  \HH^2(\Ku(X))=H^0(T_Q)
 \]
 である。cubicのmarked変形関手はsmoothかつset-valuedである。
 補題\ref{lem:induction}を適用するため、幾何学的な変形の像の自己同型を調べる。

 \(\Ku(X_R)\) の恒等に還元される自己同値を表すinvertible bimoduleは、
 \(\HH^1(\Ku(X))=0\) により相対対角bimoduleと同型である。
 これは補題\ref{lem:automorphisms}と同じsmall extensionの帰納法で従う。
 特殊ファイバー上の同型を指定しても、そこでのscalarは \(R^\times\) へ持ち上がるので、
 同型をmarkingと両立するように選べる。
 従ってmarkedな自己同値の同型類は恒等だけであり、
 必要な自己同型持ち上げ条件は満たされる。
 \(Q\) 側の自己同型についても、補題\ref{lem:theta}の比較写像が
 \(\pi_1\) 上の同型を誘導するため同じことが成り立つ。

 補題\ref{lem:induction}から、cubic変形と \(\pi_0Q(R)\) の全単射が得られる。
 幾何学的変形の像はactual smooth proper変形であり、
 補題\ref{lem:theta}によりactual変形間の同値を検出できる。
 よって最初の主張が従う。この議論はHochschild形式変形の一般的なactual実現性を
 仮定せず、この場合には幾何学的なliftで全てを実現している。

 得られた同値 \(\Ku(X_R)\simeq\cK_R\) を通じて
 \(\rho_{X_R}\) を移せば、\(\rho_X\) のliftが存在する。
 二つのmarkedなlift \(\rho_R,\rho'_R\) があれば、
 \(\rho'_R\rho_R^{-1}\) は恒等に還元されるinvertible bimoduleである。
 再び \(\HH^1=0\) による対角のrigidityを用いると、
 \(\rho'_R\rho_R^{-1}\simeq\mathrm{id}\) となり、markingを保つ同型も選べる。
 これがliftの同型類の一意性とforgetful mapの単射性を与える。
 各small extensionで還元上の同値を保って構成できるため、
 形式族での互換性も従う。
\end{proof}

\begin{remark}\label{rem:unpolarized-scope}
 この定理の一意性は同型類についてのものである。
 rotationとのintertwinerにはscalar自由度があり、
 rotationを忘れる写像がfull higher deformation groupoidの同値だとは主張しない。
 実際、\(\HH^0=\CC\)、\(\HH^1=0\) とrotationのscalar上の自明作用から、
 pairの接複体には \(H^{-1}(T_P)\simeq\CC\) がある一方、
 \(H^{-1}(T_Q)=0\) である。

 \(n=4\) では幾何学的な20方向に加えて \(\HH^2\) に二方向があり、
 rotationの固定部分を取って初めて接写像が同型となる。
 この次元でrotationを省く結論を上の証明から導くことはできない。
 ここでは、Hochschildの全22方向が各Artin底上でactual変形として
 実現されるとは主張していない。

 \(n=7\) では \(\HH^3\simeq\CC^2\) が残るが、
 幾何学的なliftと既に存在するliftを比較する帰納法により結論は成り立つ。
 障害を収め得るcohomology群が非零であることと、
 この変形問題で実際に障害が現れることは同じではない。
 一般の有限パラメータの圏からrotationを省けるという大域的主張も含めない。
\end{remark}

\section{飽和な代数的族に対するétale局所復元}\label{sec:algebraization}

形式的な復元を有限パラメータへ移す際には、圏の族そのものの代数性が必要である。
本節ではそれを明示的に仮定し、可換モデルと圏同値の両方をétale近傍上で構成する。
各次数がflatであることだけから、この仮定が自動的に従うとは主張しない。

\begin{definition}\label{def:algebraic-family}
 \(S\) を有限型複素スキームとする。
 本節で\emph{smooth properな代数的dg圏の族}とは、
 étale局所的に \(\Perf(A)\) と書け、\(A\) が底環上のdg代数で、
 \(A\) 自身が底上perfect、対角 \(A\) が \(A^{\mathrm{op}}\otL A\) 上perfectである族をいう。
 圏は小さく、pretriangulatedかつ冪等完備とし、すべての底変更を導来的に取る。
 自己同値 \(\rho\) の代数的族は、可逆なperfect bimoduleによって与える。
\end{definition}

\begin{theorem}[étale局所的な可換復元]\label{thm:etale}
 \(S\) を有限型複素スキーム、\(s\in S(\CC)\) とし、
 \((\cK_S,\rho_S)\) を定義\ref{def:algebraic-family}の組とする。
 滑らかなcubic \(X^n\subset\PP^{n+1}\)、\(n\ge3\)、およびenhanced pair同値
 \[
 \phi_s:(\cK_s,\rho_s)\xrightarrow{\sim}(\Ku(X),\rho_X)
 \]
 が与えられたとする。このとき、あるpointed étale近傍
 \((T,t)\to(S,s)\)、滑らかな相対cubic \(Y\subset\PP^{n+1}_T\)、および
 \[
 \Phi_T:(\cK_T,\rho_T)\xrightarrow{\sim}(\Ku(Y/T),\rho_{Y/T})
 \]
 が存在し、\(t\) 上で \(Y_t=X\)、\(\Phi_t\simeq\phi_s\) とできる。
 この同値はperfect bimoduleとrotationを結ぶ同型によって与えられる。

 従って、可換cubicのpairで実現できる複素点は \(S\) のZariski開部分をなす。
 その各複素点で、実現する滑らかなcubicの射影同型類は一意である。
\end{theorem}

この定理は、底 \(S\) 全体上のcubicの一意な降下や、同値の一意性を述べない。
また、\(S\) が連結であることだけでは、上の開部分が \(S\) 全体であるとはいえない。
\(n\ne4\) ではrotationを最初から与えなくても、同じ方法でétale局所的な
可換モデルが存在する（系\ref{cor:etale-unpolarized}）。

\subsection{指定された還元同値を保つ持ち上げ}

\begin{lemma}\label{lem:relative-realization}
 \(R'\twoheadrightarrow R\) を剰余体 \(\CC\) の局所Artin代数のsmall extensionとする。
 \((\cK_{R'},\rho_{R'})\) をsmooth proper pairとし、滑らかなcubic \(Y_R\) と同値
 \[
 \phi_R:(\cK_R,\rho_R)\simeq(\Ku(Y_R),\rho_{Y_R})
 \]
 を固定する。このとき \(Y_R\) を延長する \(Y_{R'}\) と、
 還元後に \(\phi_R\) と同型になるpair同値 \(\phi_{R'}\) が存在する。
 \(Y_R\) の射影空間への埋込みも指定してよい。
\end{lemma}
\begin{proof}
 特殊ファイバーでの \(\phi_R\) をmarkingとして固定する。
 まずcubicの係数を持ち上げ、幾何側に一つのlift \(Y^0_{R'}\) を取る。
 \(P\) を第\ref{sec:formal}節のpairのformal moduli problemとし、
 \(I=\ker(R'\to R)\) と置く。
 与えられた圏とこの幾何学的liftを、\(\phi_R\) によって
 \(P(R')\to P(R)\) の同じhomotopy fibre内で比較する。
 存在するliftの差は \(H^0(T_P)\otimes I\) のtorsorで測られる。
 幾何学的Kodaira--Spencer写像
 \(H^1(T_{Y_0})\simeq H^0(T_P)\) により \(Y^0_{R'}\) を修正し、その差を消す。
 補題\ref{lem:theta}の完全忠実性で、指定した還元同値を保つactual同値へ戻る。
 途中で還元上の自己同型が現れる表示を用いた場合には、
 補題\ref{lem:automorphisms}でそれを持ち上げて補正できる。

 抽象的な偏極cubicのliftに対しても、\(\OO(1)\) の延長と
 \(H^0(\OO(1))\) の自由な基底の持ち上げにより、指定された埋込みを延長できる。
 これは \(\operatorname{PGL}_{n+2}\) のsmoothnessを用いることと同じである。
\end{proof}

\subsection{同値とrotationの同型を表す代数的スタック}

\(U\subset\PP(\CC[x_1,\ldots,x_{n+2}]_3)\) を滑らかなcubicの方程式空間とし、
 \(\mathcal X\to U\) を普遍cubicとする。
 相対的な例外列とmutationにより
 \((\Ku(\mathcal X/U),\rho_U)\) はsmooth properな代数的pairの族となる。
 相対生成子として \(\bigoplus_{i=0}^{n+1}\OO_{\mathcal X}(i)\) の残余射影を使える。
 この生成性は、射影空間の相対Beilinson生成子の引戻しと、
 閉埋込みの押出しが保守的であることから従う。

\begin{proposition}\label{prop:equivalence-stack}
 \(S\)-スキーム \(T\) に対し、
 \[
 (f:T\to U,\ P,\ \alpha:P\circ\rho_T\xrightarrow{\sim}\rho_f\circ P)
 \]
 を対象とするスタック \(\mathcal E\) を考える。
 ここで \(P:\cK_T\simeq\Ku(\mathcal X_f/T)\) はperfect bimoduleで表される同値、
 射は \(\alpha\) と両立するbimoduleの同型である。
 このスタックは \(S\times U\) 上、従って \(S\) 上で局所有限表示のArtin 1-stackである。
\end{proposition}
\begin{proof}
 すべての構成はétale降下と両立するため、底をaffineにして生成子を選んでよい。
 \(\cD=\cK_S^{\mathrm{op}}\otL\Ku(\mathcal X/U)\) は \(S\times U\) 上smooth properである。
 Toën--VaquiéのTheorem 3.6とCorollary 3.15 \cite{TV}により、
 そのperfect bimoduleのmoduliは局所幾何的で局所有限表示である。
 同論文は一般の可換底環上で定式化されているので、定数族だけに限定されない。

 \(P\) が同値を与える条件は開である。
 実際、右随伴を表すperfect bimoduleを \(P^R\) とすると、
 単位と余単位
 \[
 \Delta_{\cK}\longrightarrow P^R P,
 \qquad PP^R\longrightarrow\Delta_{\Ku}
 \]
 のconeがともに零となることが必要十分である。
 生成子でdg代数を選んだ模型で、これらのconeを評価する。
 smooth proper性によりその底上の複体はperfectであり、消滅はそこで検出される。
 従って消滅座は開である。
 これは \cite[Lemma 3.25とCorollary 3.24の証明]{TV} の二つの圏に対する同じ議論である。

 同値 \(P\) の幾何ファイバーで
 \[
 \Ext^i(P,P)=\HH^i(\Ku(X_f)),\qquad i\in\ZZ.
 \]
 補題\ref{lem:HH}から負次数は消える。
 perfectな自己Hom複体のTor振幅の開性により、これは任意の底変更後にも負次数が
 消える開部分として扱える。\cite[Corollary 3.21]{TV}により、
 同値を表す部分は通常のArtin 1-stackとなる。

 最後に \(P\rho_T\) と \(\rho_fP\) の間の同型を付加する。
 射のmoduliの局所有限表示性は \cite[Corollary 3.19]{TV} による。
 その中でconeが零となる部分を取れば同型のmoduliを得る。
 同型が存在する点では両側は可逆bimoduleであり、
 負次数のHomは再び \(\HH^{<0}\) の消滅によって零である。
 従ってこの同型の選択は代数的空間で表され、
 \(\mathcal E\) はArtin 1-stackのままである。
 この構成には \(\rho^3\simeq[2]\) の選択や、そのcoherenceは課していない。
\end{proof}

\subsection{smoothnessとétale局所section}

\begin{proof}[定理\ref{thm:etale}の証明]
 \(\phi_s\) と \(X\) の方程式を選んで \(e\in\mathcal E(\CC)\) を得る。
 \(V\to\mathcal E\) をsmooth atlasとし、\(e\) の上に複素点 \(v\) を選ぶ。
 \(v\) における \(V\to S\) のsmall extensionの持ち上げ問題を考える。
 補題\ref{lem:relative-realization}により、対応する \(\mathcal E\) の対象を
 指定された還元を保って持ち上げられる。
 さらに \(V\to\mathcal E\) のsmoothnessにより、それを \(V\) へ持ち上げる。
 底が局所Artin環で剰余体が \(\CC\) なので、atlasへのétale局所持ち上げはsectionを持ち、
 この持ち上げを同じ底上で取れる。

 \(V\to S\) は局所有限型であるから、Noetherian底上の無限小判定
 \cite[Tag 02HX]{Stacks}により \(v\) でsmoothである。
 \(v\) の近傍を縮め、smooth morphismのétale局所section
 \cite[Tag 055Uとその証明中のLemma 37.38.5]{Stacks}を使うと、
 \(\kappa(v)=\kappa(s)=\CC\) であることから、指定した \(v\) を通る
 \((T,t)\to(S,s)\) と \(T\to V\to\mathcal E\) が得られる。
 これが求める \(Y\)、\(P\)、\(\alpha\) を同時に与える。

 \(T\to S\) の像はZariski開であり、その各複素点の上には複素点を取れる。
 従って可換モデルが存在する複素点の集合は開部分の複素点集合である。
 各点でのモデルの一意性は \cite[Theorem 1.7]{LZ} による。
\end{proof}

\begin{corollary}\label{cor:generic-realization}
 定理\ref{thm:etale}の仮定で \(S\) が既約なら、\(s\) を含む空でないZariski開部分
 \(S^\circ\) の全複素ファイバーに可換cubicモデルが存在する。
 幾何学的generic fibreにも、有限型étale底変更を通じてモデルが存在する。
\end{corollary}
\begin{proof}
 上のétale近傍の像を \(S^\circ\) とする。既約性からこれはdenseである。
 \(S\) を必要なら被約化してgeneric pointを考えると、
 \(T\) の支配的な成分のgeneric residue fieldは \(\CC(S)\) の有限分離拡大である。
 その上のcubicと同値を代数閉包へ底変更すればよい。
\end{proof}

\begin{remark}\label{rem:algebraization-boundary}
 本節では形式的なpower seriesの収束を証明していない。
 指定された代数的圏の族に対し、代数的な同値のmoduliのsmoothnessから
 別途、実際の代数的近傍と同値を作っている。
 その完備化は形式的復元定理のモデルとmarked同型類で一致する。
 しかし、任意に選ばれた形式的座標表示や形式的bimoduleそのものが
 元の底上へ一意に降下する、という結論ではない。
\end{remark}


\begin{corollary}[rotationを指定しない局所的存在]\label{cor:etale-unpolarized}
 定理\ref{thm:etale}で \(n\ne4\) とする。
 \(\rho_S\) は仮定せず、\(\cK_s\simeq\Ku(X)\) というenhanced同値だけを与えても、
 pointed étale近傍 \(T\) 上にcubic \(Y/T\) と、指定された特殊ファイバーの同値を
 延長する \(\cK_T\simeq\Ku(Y/T)\) が存在する。
\end{corollary}
\begin{proof}
 命題\ref{prop:equivalence-stack}から \(\alpha\) の指定を除いた同値のスタックを用いる。
 定理\ref{thm:unpolarized-formal}の証明を、還元同値を指定したhomotopy fibre内で行えば、
 補題\ref{lem:relative-realization}のrotationなしの版が得られる。
 後は定理\ref{thm:etale}と同じsmooth atlasとétale局所sectionの議論を適用する。
 得られたモデルからrotationを移せるが、その \(T\) 全体上での一意性は主張しない。
\end{proof}

\section{Sklyaninブロックと次数シフト}\label{sec:elliptic}

\subsection{中心元と点曲線}

\begin{lemma}\label{lem:center}
 \(G_h=x^3+y^3+z^3-2h^2\sym(xyz)\) は、すべての \(h\) で
 \(S_h\) の非零な中心正則元である。
\end{lemma}
\begin{proof}
 \(D=x^3+y^3+z^3\) とする。関係式から
 \[
 [D,x]=h(zy^2z-yz^2y)=h[zy,yz]=h^2[zy,x^2].
 \]
 また六つの順列を同じ順序へ移すと
 \[
 \sym(xyz)=xyz+\frac h2(-x^3+y^3-z^3)
 \]
 となる。この式を \(x\) と交換させ、再び二次関係式を使えば
 \(2[\sym(xyz),x]=[zy,x^2]\) を得る。
 従って \([G_h,x]=0\)。巡回置換で \(y,z\) についても従う。

 非零性は特殊化で失われない形で確認できる。
 自由代数の次数三部分の線形汎関数 \(\ell\) を
 \[
 \ell(x^3)=\ell(y^3)=\ell(z^3)=1,\quad
 \ell(xyz)=\ell(yzx)=\ell(zxy)=h/2,
 \]
 \[
 \ell(xzy)=\ell(zyx)=\ell(yxz)=-h/2
 \]
 とし、ちょうど二文字が等しい語では零とする。
 三つの二次関係式の左右に一文字を掛けた18個の式を \(\ell\) はすべて殺す。
 例えば \(\ell(z(xy-yx-hz^2))=h/2+h/2-h=0\) である。
 従って \(\ell\) は \((S_h)_3\) に降り、\(\ell(G_h)=3\) となる。
 \((a,b,c)=(1,-1,-h)\) はSklyanin代数の12個の非正則パラメータに属さず、
 \(S_h\) はAS-regular domainである \cite[Introduction]{CGWZ}。
 \(h=0\) は多項式環である。よって \(G_h\) は左右正則である。
\end{proof}

標準Sklyanin記法では \((a,b,c)=(1,-1,-h)\) である。
標準中心元の式 \cite[Equation (2.16)]{Walton}
\[
 g=c(a^3-c^3)x^3+a(b^3-c^3)xyz+b(c^3-a^3)yxz+c(c^3-b^3)y^3
\]
に代入し、\(yxz=xyz-hz^3\) を用いると \(g=-hG_h\) となる。
\(h\ne0\)、\(h^6\ne-27\) では点曲線の公式
\cite[Theorem 2.14]{Walton}から
\[
 E_h=V(x^3+y^3+z^3+h^2xyz)\subset\PP^2
\]
を得る。これは滑らかな楕円曲線であり、
\[
 S_h/(G_h)\simeq B(E_h,L,\sigma_h),\qquad L=\OO_{E_h}(1),
\]
と書ける。ここで \(\sigma_h=t_p\) は平行移動である。
原点をinflection point \(o=[1:-1:0]\) に取ると、
多重線形化した関係式を第一成分 \(o\) で評価した核は
\(p=[1:-1:-h]\) である。従ってこの規約で \(\sigma_h=t_p\) となる。
この中心商と楕円曲線の比較は \cite[Sections 3.5--3.6]{EG} も参照。
\(h=0\) は通常のFermat平面三次曲線の斉次座標環である。
この場合は \(p=o\)、\(\sigma_0=\Id\) とする。

\subsection{rotationに含まれる平行移動を消す}

\begin{lemma}\label{lem:elliptic-conjugacy}
 \(E\) を滑らかな楕円曲線、\(L\) を次数三の直線束、
 \(\sigma=t_p\) とする。
 \(C=B(E,L,\sigma)\) のgraded singularity categoryを \(D^b(E)\) と同一視したとき、
 その次数シフトは可換平面三次曲線のrotation
 \(\rho_0=T_{\OO_E}\circ(-\otimes L)\) と共役である。
\end{lemma}
\begin{proof}
 右加群の規約を
 \[
 \Gamma_*(F)_d=H^0(E,F\otimes L_d),\qquad
 L_d=L\otimes\sigma^*L\otimes\cdots\otimes(\sigma^{d-1})^*L
 \]
 と取る。tailsの次数シフトは \(H(F)=\sigma_*(F\otimes L)\) である。
 Gorenstein parameter零のOrlov windowを比較すると
 \[
 \rho_{\mathrm{nc}}=T_{\OO_E}\circ H=t_{-p}^*\rho_0.
 \tag{2}\label{eq:rho-nc}
 \]
 このwindow公式を確認する。特異圏への商を \(q\) とし
 \(\Psi_i=qR\Gamma_{\ge i}\) と置くと、
 \((1)\Psi_1=\Psi_0H\)、\(q(k)=\Psi_0(\OO_E)[1]\)、
 \[
 \Psi_1F\longrightarrow\Psi_0F
 \longrightarrow\RHom(\OO_E,F)\otimes q(k)
 \]
 がある。楕円曲線のSerre双対から
 \(\Psi_1=\Psi_0T_{\OO_E}^{-1}\) となり、\eqref{eq:rho-nc}を得る。

 原点を \(o\)、\(P_q=\OO_E(q-o)\)、\(A_q=(-\otimes P_q)\) とする。
 spherical twistの核は \(I_\Delta[1]\) なので、
 \[
 K_{T_{P_q}}=(P_q^{-1}\boxtimes P_q)\otimes I_\Delta[1].
 \]
 see-sawを両方向のファイバーへ適用すると
 \[
 (P_q^{-1}\boxtimes P_q)\otimes I_\Delta
 \simeq I_{\{y=x-q\}}.
 \]
 よって \(T_{P_q}=t_q^*T_{\OO_E}\) であり、
 \[
 A_q\rho_0A_q^{-1}=t_q^*\rho_0,\qquad
 \rho_{\mathrm{nc}}=A_{-p}\rho_0A_{-p}^{-1}.
 \tag{3}\label{eq:conjugacy}
 \]
 これらは核とwindowによるdg同値として実現される。
 三次周期性の二つの同型の差は \(\HH^0(D^b(E))^\times=\CC^\times\) のscalarであり、
 共役の自然同型をその三乗根で調整すれば周期性も両立させられる。
 \(\HH^{<0}(D^b(E))=0\) なので、この自然同型の等式に追加の高次ホモトピー障害はない。
 後で用いる有限次数差の作用は、この周期性から構成する。
\end{proof}

\begin{corollary}\label{cor:block}
 \(h^6\ne-27\) のとき、三次周期性と次数シフトを保つdg Morita同値
 \[
 \MF^{\ZZ}(S_h,G_h)
 \simeq
 \MF^{\ZZ}\bigl(\CC[x,y,z],x^3+y^3+z^3+h^2xyz\bigr)
 \]
 が存在する。
\end{corollary}
\begin{proof}
 \cite[Theorem 3.6.2]{EG}、Orlov window、
 \cite[Theorem 5.8]{CCKM}によって両辺を \(D^b(E_h)\) と比較し、
 補題\ref{lem:elliptic-conjugacy}を使う。
 MFと特異圏の比較は、factorizationを商環上の両側無限の完全周期分解へ送る
 dg関手で実現する。中心正則元の商はGorensteinであり、complete resolutionは
 MCM加群の安定圏のdg模型を与える。
 CCKMの完全忠実性をすべてのshiftに適用すると、Hom複体の比較がquasi-isomorphismとなる。
 周囲がAS-regularであることによる高次syzygyのfactorization表示が本質的全射性を与える。
 特異圏からwindowへのdg比較は \cite[Section 5, Theorem 5.8]{CT} による。
 \(h=0\) では恒等の可換比較を用いる。
\end{proof}

\section{中心三次式を加える比較と通常次数への降下}\label{sec:TS}

\subsection{二つの次数ブロックと外部積}

中心変数側の三次式をFermat型に制限する必要はない。
この節では \(T=\CC[u_1,\ldots,u_m]\)、\(m\ge1\)、とし、
\(P\in T_3\) は
\[
 J(P)=\left(\frac{\partial P}{\partial u_1},\ldots,
                 \frac{\partial P}{\partial u_m}\right),\qquad
 \dim_\CC T/J(P)<\infty                                      \tag{4a}\label{eq:central-isolated}
\]
を満たすものとする。これは斉次三次式 \(P\) のアフィン臨界点が
原点だけであるという条件である。
\[
 \Gamma=(\ZZ e_0\oplus\ZZ e_1)/\langle3e_0-3e_1\rangle,
 \qquad w=3e_0=3e_1                                       \tag{4b}\label{eq:central-gamma}
\]
と置く。\(S_d\) の次数を \(de_0\)、すべての \(u_i\) の次数を \(e_1\) とする。
したがって中心変数どうしの混合項を含む \(P\) も次数 \(w\) である。

\begin{lemma}[任意の孤立中心三次式に対する外部積]\label{lem:TS}
 \(S\) を連結非負次数付きnoetherian AS-regular \(\CC\)-代数、
 \(g\in S_3\) を中心正則元とする。
 \eqref{eq:central-isolated}の下で、外部積はpretriangulated hullと
 冪等完備化の後にdg Morita同値
 \[
 \MF^\ZZ(S,g)\otL_\CC\MF^\ZZ(T,P)
 \simeq
 \MF^\Gamma(S\otimes_\CC T,g+P)                             \tag{4c}\label{eq:central-TS}
 \]
 を与える。この同値は二つの因子の次数シフトと、
 両者の三次内部シフトを \([2]\) と同一視する周期性を保つ。
\end{lemma}
\begin{proof}
 \(R=S\otimes_\CC T=S[u_1,\ldots,u_m]\)、\(W=g+P\)、
 \(A=R/(W)\)、\(C=S/(g)\) と置く。
 まず \(W\) は中心正則元である。
 実際、\(u\) に関する単項式順序を固定すると、\(P\ne0\) の最高項の係数は
 非零のscalarであり、任意の非零多項式との積の最高項は消えない。
 以下では中心正則元の商に対するMFとgraded singularity categoryの比較
 \cite[Theorem 5.8]{CCKM}を用いる。
 \(\Gamma\to\ZZ\)、\(e_0,e_1\mapsto1\)、は有限の核を持ち、
 正の整数次数を与えるので、有限次数付き自由分解による同じ比較は
 \(\Gamma\)-次数についても適用できる。

 \paragraph{外部積と完全忠実性.}
 factorizationの規約を
 \[
 Q=(Q_1\xrightarrow{\varphi}Q_0
                  \xrightarrow{\psi}Q_1(w)),\qquad
 \psi\varphi=g\Id,\quad\varphi\psi=g\Id
 \]
 とする。これを \(Q^{i+2}=Q^i(w)\) を満たす曲がった複体に展開すると、
 \(Q^0=Q_0\)、\(Q^1=Q_1(w)\) である。
 二つのfactorizationの外部積は、二つの奇微分をKoszul符号付きで足すことにより定める。
 奇微分の外部和の平方は \(g+P\) である。
 二つの \([2]\) を同じ内部シフトに対応させる関係が
 \(3e_0=3e_1\) であり、得られるfactorizationは有限自由である。

 有限自由対象 \(Q,Q'\) と \(K,K'\) に対して、
 係数を \(S\)-次数と \(T\)-次数に分解すると
 \[
 \Hom_{\MF^\Gamma}(Q\boxtimes K,Q'\boxtimes K')
 \simeq
 \operatorname{Tot}\!\left(
   \Hom_{\MF^\ZZ(S,g)}(Q,Q')\otimes_\CC
   \Hom_{\MF^\ZZ(T,P)}(K,K')\right)                         \tag{4d}\label{eq:central-hom}
 \]
 を得る。具体的には、左辺の次数 \(r\) の斉次写像は、
 \(a+b=r\) を満たす右辺の次数 \((a,b)\) の写像の和に一意に分解する。
 二つの内部次数の差を移す操作は \((3,-3)\) の整数倍だけであり、
 これは二つのcohomological degreeを \((2,-2)\) だけ変える操作に対応する。
 微分と合成の符号も外部和のKoszul符号に一致する。
 各因子のHom複体は、有限個の自由加群の次数シフトと係数代数の非負次数性から
 下に有界である。従って各全cohomological degreeに寄与する項は有限個であり、
 \eqref{eq:central-hom}に完備化は必要ない。
 これが外部積のdg完全忠実性を与える。

 \paragraph{中心変数を零とした商への生成性の帰着.}
 \(\MF^\Gamma(R,W)\) の有限自由対象の奇微分を \(D\) とする。
 \(u_i\) は中心変数なので、係数行列の \(u_i\)-微分は \(R\) 上の導分であり、
 \[
 D\,\frac{\partial D}{\partial u_i}
 +\frac{\partial D}{\partial u_i}\,D
 =\frac{\partial P}{\partial u_i}\Id
 \]
 が成り立つ。内部シフトを付けると、これは
 \(\partial_iP:N\to N(2e_1)\) の零ホモトピーである。
 斉次イデアル \(J(P)\) は \((u_1,\ldots,u_m)\)-primaryなので、
 ある整数 \(d\ge2\) に対してすべての \(u_i^d\) が \(J(P)\) に属する。
 斉次係数でこれらを偏微分の線形結合に書けば、
 \(u_i^d:N\to N(de_1)\) も特異圏で零となる。

 \(N\) を有限生成 \(A\)-加群の有界複体で表し、
 \(u_1^d,\ldots,u_m^d\) に沿うKoszul複体を \(K(u^d;N)\) と書く。
 零写像のconeは元の対象を直和因子に持つため、\(N\) は
 \(K(u^d;N)\) の直和因子である。
 より正確には、一変数ごとのKoszul cone関手は完全関手でretractを保つので、
 各 \(u_i^d\) の零作用から得られるretractを順に合成すればよい。
 このKoszul複体のcohomologyは \((u_1^d,\ldots,u_m^d)\) に殺され、
 \[
 (u_1,\ldots,u_m)^{m(d-1)+1}\subset(u_1^d,\ldots,u_m^d)
 \]
 である。従って各cohomology加群は、有限個の \(C=A/(u_1,\ldots,u_m)\)-加群を
 商とする有限filtrationを持つ。標準切断の三角形も使うと、
 右辺の特異圏は \(C\)-加群の制限とその次数シフトで厚く生成される。

 ここで制限が特異圏へ降りることを確認する。
 \(u_1,\ldots,u_m\) は \(A\) の中心正則列である。
 一つの変数 \(u_i\) について、\(R\) で \(u_i a=Wb\) とする。
 \(u_i=0\) とした商では、\(\overline W=g+P|_{u_i=0}\) は正則である。
 特殊化した \(P\) が非零なら最高項のscalar係数を使い、
 零なら \(g\) の正則性を使えばよい。
 従って \(\overline b=0\)、すなわち \(b=u_i b_1\) であり、
 多項式環 \(R\) で \(u_i\) を消去して \(a=Wb_1\) を得る。
 先行する変数を順に零とした同じ議論が正則列を与える。
 よって \(C\) は有限自由Koszul分解を持つ \(A\)-perfect加群であり、
 制限は \(\Dsg(C)\to\Dsg(A)\) を誘導する。

 \paragraph{制限加群の安定化の明示.}
 Euler公式から
 \[
 q_i=\frac13\frac{\partial P}{\partial u_i},\qquad
 P=\sum_{i=1}^{m}u_iq_i
 \]
 と書ける。\(T\) 上の通常のKoszul分解を
 \[
 K_j=T(-je_1)\otimes_\CC\Lambda^j\CC^m,
 \qquad d_K=\sum_i u_i\iota_i
 \]
 とする。\(\iota_i\) は外積代数の縮約、\(\varepsilon_i\) は対応する基底による左外積である。
 \[
 h_K=\sum_iq_i\varepsilon_i:K_j\longrightarrow K_{j+1}(w)
 \]
 と置くと
 \[
 d_Kh_K+h_Kd_K=P\Id,\qquad h_K^2=0.                         \tag{4e}\label{eq:central-koszul}
 \]
 ここで \(q_i\) の次数は \(2e_1\) であり、
 \(K_{j+1}(w)\) のシフトによりすべての表示写像は内部次数零である。

 有限Koszul複体を偶奇で折り畳み、
 \[
 (K_P)_0=\bigoplus_{j\ge0}K_{2j}(jw),\qquad
 (K_P)_1=\bigoplus_{j\ge0}K_{2j+1}(jw)
 \]
 と置く。\(d_K+h_K\) は
 \[
 (K_P)_1\longrightarrow(K_P)_0
             \longrightarrow(K_P)_1(w)
 \]
 を与え、\eqref{eq:central-koszul}によってこれは \(P\) の有限自由factorizationである。

 \(M\) を \(C\)-加群とし、その特異圏での像を表すfactorizationを
 \(Q=(Q_1\xrightarrow\varphi Q_0\xrightarrow\psi Q_1(w))\) とする。
 有限回のsyzygyとcohomological shiftを取れば
 \(M=\operatorname{coker}\varphi\) としてよい。
 実際、元の加群との差に現れる \(C\)-perfect対象は制限後も \(A\)-perfectである。
 \(\psi\varphi=g\Id\) と \(g\) の正則性から \(\varphi\) は単射であり、
 \[
 0\longrightarrow Q_1\xrightarrow\varphi Q_0\longrightarrow M\longrightarrow0
 \]
 は \(S\)-自由分解である。
 これをhomological degree \(1,0\) に置き、
 \(d_Q=\varphi\)、\(s_Q:Q_0\to Q_1(w)\) を \(\psi\) とする。
 \(s_Q\) は \(Q_1\) 上では零とする。

 \(u_iM=0\) として \(M\) を \(R\)-加群に制限したときの有限自由分解は
 \[
 F_a=\bigoplus_{b+c=a}Q_b\otimes_\CC K_c,
 \qquad d=d_Q\otimes1+(-1)^b\otimes d_K
 \]
 である。\(b\in\{0,1\}\) は \(Q\) 側のhomological degreeである。
 \[
 s=s_Q\otimes1+(-1)^b\otimes h_K:F_a\longrightarrow F_{a+1}(w)
 \]
 と置くと、異なる因子の奇作用の反可換性から
 \[
 ds+sd=W\Id,\qquad s^2=0.                                  \tag{4f}\label{eq:central-shamash}
 \]
 が成り立つ。

 このデータの周期化が実際に \(A\)-自由分解を与えることも確認しておく。
 \(F_a=0\) for \(a<0\) とし、
 \[
 \mathcal F_n=
 \bigoplus_{j\ge0}(F_{n-2j}\otimes_R A)(-jw),\qquad n\ge0
 \]
 と置く。微分は同じ \(j\) 内の \(d\) と、\(j\ge1\) について
 \[
 s:(F_{n-2j}\otimes_R A)(-jw)
 \longrightarrow(F_{n-2j+1}\otimes_R A)(-(j-1)w)
 \]
 の和である。\eqref{eq:central-shamash}により平方は零になる。
 各 \(n\) で和は有限である。

 \(j\) による増大濾過を使うと、第一ページは
 \(H_*(F\otimes_R A)\) の内部・homological shiftの和となる。
 \(W\) が正則で \(WM=0\) なので、
 \[
 H_i(F\otimes_R A)=
 \begin{cases}
 M&i=0,\\
 M(-w)&i=1,\\
 0&i\ne0,1.
 \end{cases}
 \]
 次の微分は、\(j\ge1\) の \(H_0\) を \(j-1\) の \(H_1\) へ送る同型である。
 これを \(ds+sd=W\Id\) から直接見られる。
 \(a\in F_0\) が \(M\) の元を表すなら、\(s(a)\) の \(A\) 上の像はcycleであり、
 \(d s(a)=Wa\) である。
 短完全列
 \[
 0\longrightarrow F(-w)\xrightarrow{W}F
      \longrightarrow F\otimes_R A\longrightarrow0
 \]
 のconnecting mapはそのcycleを元の \(a\) に送る。
 従ってこの微分は、上の \(H_1\simeq M(-w)\) の標準同定の下で恒等である。
 隣接する二項はすべて相殺し、\(j=0\) の \(H_0=M\) だけが残る。
 濾過は各全homological degreeで有限なので収束し、
 \(\mathcal F\) は \(M\) の \(A\)-自由分解となる。

 \(F\) は有限複体なので、この分解の十分先は二周期的である。
 そのfactorizationは
 \[
 \widehat F_0=\bigoplus_{j\ge0}F_{2j}(jw),\qquad
 \widehat F_1=\bigoplus_{j\ge0}F_{2j+1}(jw),
 \qquad d+s
 \]
 による折り畳みであり、Koszul符号と内部シフトを含めて
 \(Q\boxtimes K_P\) に一致する。
 実際、\(\widehat F_0\) は
 \(Q_0\otimes(K_P)_0\oplus Q_1\otimes(K_P)_1(w)\)、
 \(\widehat F_1\) は
 \(Q_1\otimes(K_P)_0\oplus Q_0\otimes(K_P)_1\) である。
 従って制限加群の安定化は
 \[
 \operatorname{stab}_A(M)\simeq Q\boxtimes K_P.               \tag{4g}\label{eq:restriction-mf}
 \]
 ここで全体のshiftにも曖昧さはない。
 十分大きい \(\ell\) について
 \(\mathcal F_{2\ell}=(\widehat F_0\otimes_R A)(-\ell w)\)、
 \(\mathcal F_{2\ell+1}=(\widehat F_1\otimes_R A)(-\ell w)\) であり、
 そのcokernelは \(\Omega_A^{2\ell}M\) である。
 \(\Omega_A^{2\ell}M\simeq M[-2\ell]\) と \((w)\simeq[2]\) によって
 この内部シフトはちょうど打ち消される。

 最後に、\(\Gamma/\ZZ e_0\simeq\ZZ/3\) であるため、
 任意の \(\Gamma\)-次数付き \(C\)-加群は、
 三つの次数類 \(0,e_1,2e_1\) に属する部分に分解する。
 各部分は通常の \(\ZZ\)-次数付き \(C\)-加群の内部シフトであり、
 \eqref{eq:restriction-mf}の右辺の \(T\) 側を \((j)\)、\(j=0,1,2\)、
 だけシフトすればすべて得られる。
 前段のKoszul filtrationと合わせて、外部積の像は右辺全体を厚く生成する。
 これで本質的全射性が従う。
 周期性および二つの次数シフトとの両立は、ここで用いた外部和の構成そのものから従う。
\end{proof}

\begin{remark}[一変数の場合の明示行列]\label{rem:central-one-variable}
 \(m=1\)、\(P=u^3\) では \(K_P=(u,u^2)\) である。
 上の分解とホモトピーは
 \[
 Q_1(-e_1)\xrightarrow{(-u,\varphi)^t}
 Q_1\oplus Q_0(-e_1)\xrightarrow{(\varphi,u)}Q_0,
 \]
 \[
 s_0=(\psi,u^2)^t,\qquad s_1=(-u^2,\psi)
 \]
 となる。折り畳んだfactorizationは
 \[
 \alpha=
 \begin{pmatrix}\varphi&u\\-u^2&\psi\end{pmatrix}:
 Q_1\oplus Q_0(-e_1)\longrightarrow Q_0\oplus Q_1(w-e_1),
 \]
 \[
 \beta=
 \begin{pmatrix}\psi&-u\\u^2&\varphi\end{pmatrix}:
 Q_0\oplus Q_1(w-e_1)\longrightarrow
 Q_1(w)\oplus Q_0(w-e_1).
 \]
 中心性により \(\alpha\beta=\beta\alpha=(g+u^3)\Id\) である。
 これらの式は、一般の \(K_P\) による安定化の符号と次数シフトの具体的な検算になる。
\end{remark}

\subsection{有限群による整数次数の回復}

\(\Gamma\to\ZZ\)、\(e_0,e_1\mapsto1\)、の核は
\(T_\Gamma\simeq\ZZ/3\) であり、\(e_1-e_0\) で生成される。
分裂 \(1\mapsto e_0\) の下で、その双対群
\(H_\Gamma=\Hom(T_\Gamma,\CC^\times)\simeq\mu_3\) は
\(S\) 上恒等、すべての \(u_i\) を同時に三乗根倍する作用である。
\(P\) は斉次三次式なのでこの作用は \(W=g+P\) を保ち、
\[
 \MF^\Gamma(R,W)
 \simeq\bigl(\MF^\ZZ(R,W)\bigr)^{H_\Gamma}.                \tag{5}\label{eq:eq}
\]
有限群の位数は可逆である。
同変な有限次数付き自由加群は、連結な整数次数と有限群表現の半単純性により
同変な斉次基底を持つので、この比較は有限自由factorizationのdg模型で実現する。
さらに双対群 \(T_\Gamma\) で同変化すると、元の通常次数の圏が回復される
\cite[Theorem 4.2, Section 8, Lemma 8.6, Theorem 8.9]{Elagin}。
冪等完備化はこの比較の一部として取る。

外部積 \eqref{eq:central-TS} の左辺で双対作用の生成元は
\[
 \rho_T\rho_S^{-1}
\]
である。ここで \(\rho_S,\rho_T\) はそれぞれの因子の内部次数一のシフトである。
その三乗を恒等と同一視する自然同型は、両因子の
\((3)\simeq[2]\) から得られる。
二つのperiodicityを含めて \(\rho_S\) を保つdg同値はこの作用と両立し、
\eqref{eq:eq}と双対同変化により通常次数へ降りる。
この際 \(e_0\) のシフトも保存されるので、通常次数のdegree shift \((1)\) が保存される。
有限群作用のcoherenceはfactorizationの内部シフトのcoherenceから引き継ぎ、
単に自己同値の三乗が同型になるという存在命題だけから群作用を作るのではない。

\subsection{Sklyaninブロックへの適用}

\begin{proof}[定理\ref{thm:global}の証明]
 補題\ref{lem:TS}を \((S_h,G_h)\) と
 \((\CC[x,y,z],x^3+y^3+z^3+h^2xyz)\) のそれぞれに適用する。
 系\ref{cor:block}は第一因子の次数シフトと三次周期性を保ち、
 \(\MF^\ZZ(T,P)\) 側では恒等を用いる。
 前節の降下によって通常次数のdegree shiftを保つdg Morita同値
 \[
 \MF^\ZZ(S_h[u],G_h+P)
 \simeq
 \MF^\ZZ\!\left(\CC[x,y,z,u],x^3+y^3+z^3+h^2xyz+P\right)
 \]
 を得る。体上のOrlov同値と \cite{CCKM} により、これは残余圏の比較である。

 可換方程式の \(u_i\)-偏微分がすべて零なら、
 \eqref{eq:central-isolated}によって全 \(u_i=0\) である。
 \(h^6\ne-27\) の下では残る平面三次式のアフィン臨界点も原点だけである。
 よって射影超曲面 \(Y_{h,P}\) は滑らかである。
 次元は \(m+1\ge3\) なので、一意性は \cite[Theorem 1.7]{LZ} による。
\end{proof}



\begin{corollary}[平行移動を伴う一般の楕円ブロック]\label{cor:elliptic-central-cubic}
 \(S\) をHilbert series \((1-t)^{-3}\) を持つnoetherianな
 二次三次元AS-regular \(\CC\)-代数、\(g\in S_3\) を中心正則元とする。
 滑らかな平面三次曲線 \(E\)、\(L=\OO_E(1)\)、および
 ある原点に関する平行移動 \(\sigma=t_p\) に対して、次数付き同型
 \[
 S/(g)\simeq B(E,L,\sigma)
 \]
 が存在すると仮定する。
 \(f_E\in\CC[x,y,z]_3\) を \(E\subset\PP^2\) の方程式とし、
 \(P\in\CC[u_1,\ldots,u_m]_3\)、\(m\ge2\)、は
 \eqref{eq:central-isolated}を満たすとする。
 このとき \(S[u]/(g+P)\) のrotation付き残余圏の滑らかな可換モデルは
 \[
 V(f_E(x,y,z)+P(u))\subset\PP^{m+2}
 \]
 であり、射影同型を除いて一意である。
\end{corollary}
\begin{proof}
 Gorenstein parameter零のOrlov同値により
 \(\MF^\ZZ(S,g)\) は \(D^b(E)\) と同値であり、
 次数シフトは補題\ref{lem:elliptic-conjugacy}の
 \(T_{\OO_E}\sigma_*(-\otimes L)\) に対応する。
 同補題の直線束による共役とperiodicityのscalar調整により、
 \(\MF^\ZZ(\CC[x,y,z],f_E)\) との、次数シフトと三次周期性を保つdg同値を得る。
 あとは補題\ref{lem:TS}と有限次数差による降下を適用すればよい。
 滑らかさと一意性は定理\ref{thm:global}と同じ証明による。
 方程式 \(f_E\) の非零scalar倍による違いは、
 分離された \(x,y,z\) の同時三乗根倍によって射影同型に吸収される。
\end{proof}

\begin{remark}[この拡張の範囲]
 定理\ref{thm:global}では中心変数間の任意の三次混合項を許すが、
 非可換ブロックの生成元と中心変数の混合項は含めない。
 別の非可換ブロックについては、その生成元による通常の偏微分が代数の導分とは限らず、
 本節の \(J(P)\)-primaryな支持への帰着は適用できない。
 この点は独立した生成性の問題として残る。
 本節の拡張の先行研究上の新規性は未確認である。
\end{remark}

\section{複数ブロックの比較と生成性の帰着}\label{sec:multiblock}

\subsection{生成性を仮定しない外部積}

\(A_1,\ldots,A_s\) を、\((A_i)_0=\CC\) を満たす局所有限な
非負次数付き代数、\(f_i\in Z(A_i)_3\) を中心的な三次斉次元とする。
この小節の \(\MF\) は、有限自由次数付きmatrix factorizationのdg圏の
pretriangulated hullを取り、さらに冪等完備化したものを表す。
この定義にはtensor product algebraのnoetherian性を仮定しない。
\[
 \Gamma_s=\ZZ^s/\langle3e_i-3e_1:2\le i\le s\rangle,
 \qquad w=3e_1,
 \qquad A=\bigotimes_{i=1}^s A_i,
 \qquad W=\sum_{i=1}^s f_i
\]
とする。\(A_i\) の次数 \(d\) の斉次元には \(de_i\) を割り当てる。

\begin{lemma}[外部積の完全忠実性]\label{lem:multiblock-ff}
 外部積はdg完全忠実関手
 \[
 \Phi:\Perf\!\left(\bigotimes_{i=1}^s\MF^{\ZZ}(A_i,f_i)\right)
 \longrightarrow\MF^{\Gamma_s}(A,W)
 \]
 を定める。\(i\) 番目の因子のdegree shiftは右辺の \((e_i)\) に対応し、
 \((3e_i)=[2]\) の標準同型とも両立する。
\end{lemma}
\begin{proof}
 まず有限自由factorization \(P_i,Q_i\) を取る。
 tensor productの奇微分はKoszul符号付きの外部和であり、
 その二乗は \(W\) となる。
 各因子のHom複体を \(C_i^\bullet\) と書くと、
 自由加群のHomを行列係数ごとに分解することにより
 \[
 \Hom_{\MF^{\Gamma_s}(A,W)}^\bullet
       (\boxtimes_iP_i,\boxtimes_iQ_i)
 \cong
 \operatorname{Tot}\!\left(\bigotimes_i C_i^\bullet\right)
 \tag{\(\dagger\)}\label{eq:multiblock-hom}
 \]
 を得る。ここでは次数群の関係 \(3e_i=3e_1\) が、
 各因子の二つのcohomological shiftを他の因子へ移す操作に正確に対応する。
 従って右辺の次数 \(q\) の和は
 \(q_1+\cdots+q_s=q\) にわたる通常のtensor product複体の和であり、
 追加の完備化や直積を取らない。

 この和の有限性を確認する。各 \(P_i,Q_i\) は有限個の内部次数shiftからなり、
 \(A_i\) は非負次数なので、\(C_i^q=0\) は十分小さい \(q\) で成り立つ。
 したがって固定した全次数 \(q\) に寄与する \((q_1,\ldots,q_s)\) は有限個である。
 同様に、行列係数を元の \(\ZZ^s\)-次数へ分解したときも、
 固定した \(\Gamma_s\)-次数と非負次数の条件の下では有限個の成分しか現れない。
 これにより \eqref{eq:multiblock-hom} は実際のHom複体の同型である。
 differentialとcompositionの一致はtensor productのKoszul符号から従う。
 完全忠実性は有限個のcone、shift、直和因子を取っても保たれるので、主張を得る。
 内部次数shiftとperiodicityとの両立はこの構成自体に含まれる。
\end{proof}

\subsection{任意個の楕円THCRの積と周囲の正則性}

Chanの多重Serre同値 \cite[Theorem 3.4]{ChanMulti} は、代数のnoetherian性を
仮定せず、right ampleな可逆bimoduleの組に適用できる。
同論文のCorollaries 5.8--5.9は二因子のTHCRと三次元AS-regular代数の積を扱う。
以下では、その証明で用いる面と帯の議論をfactor数について帰納し、任意個の楕円因子を扱う。
Keeler \cite[Theorem 2.7, Corollary 2.8]{KeelerMulti} によっても、
この積のbimoduleの組の左右NC-amplenessを確認できる。
各translationはnumerical Picard groupに自明に作用し、全ての座標が正の外部積はampleだからである。
ただし多重amplenessだけではnoetherian性は自動的に従わず、以下で帯の有限性を別に扱う。

\begin{lemma}[楕円THCRの有限テンソル積]
\label{lem:elliptic-thcr-tensor-noeth}
 \(E_1,\ldots,E_r\) を滑らかな複素射影楕円曲線、\(L_i\) を正次数の可逆層、
 \(\sigma_i\) をtranslationとする。このとき
 \[
 C=\bigotimes_{i=1}^r C_i,\qquad C_i=B(E_i,L_i,\sigma_i)
 \]
 は左右noetherianである。
\end{lemma}
\begin{proof}
 \(X=\prod_iE_i\) とし、各factorから引き戻した可逆bimoduleを
 \(\mathcal L_i\) と書く。異なるfactorの作用は可換であり、
 \(\mathcal L^{\mathbf d}=\mathcal L_1^{d_1}\cdots\mathcal L_r^{d_r}\) と置けば
 \[
 C_{\mathbf d}=H^0(X,\mathcal L^{\mathbf d})\qquad(\mathbf d\in\mathbb N^r)
 \]
 は \(C\) のテンソル積多重次数と乗法を与える。
 可逆bimoduleの逆を用いて \(\mathcal L^{\mathbf d}\) は
 \(\mathbf d\in\mathbb Z^r\) に対しても定義する。

 \emph{Ampleness.}
 外部積のvery ample可逆層 \(H=\boxtimes_iH_i\) を固定する。
 任意のcoherent層 \(\mathcal F\) を、有限個の \(H\) の負冪の直和からの全射で
 \(\dim X\) 段まで近似する。
 各 \(\mathcal L_i^{d_i}\) の台となる可逆層の次数は \(d_i\deg L_i\) なので、
 その各近似項を \(\mathcal L^{\mathbf d}\) でtwistした高次cohomologyは、
 \(d_i\) が全て十分大きいと消える。
 ここでは楕円曲線上の正次数の可逆層の \(H^1\) の消滅とK\"unneth公式を用いる。
 最後の核にはGrothendieck消滅を適用し、dimension shiftingによって
 \[
 H^q(X,\mathcal F\mathcal L^{\mathbf d})=0
 \quad(q>0,\ \mathbf d\gg0)
 \]
 を得る。従ってこの組はright ampleである。
 \(\mathcal F\otimes H^{-q}\) にもこの消滅を適用すると、
 Castelnuovo--Mumford regularityにより
 \(\mathcal F\mathcal L^{\mathbf d}\) の大域生成も \(\mathbf d\gg0\) で従う。
 自動同型によるpushforwardはcohomologyと大域生成を変えないので、
 この議論は可逆bimoduleの左右の規約に依存しない。

 \emph{面と有限幅の帯.}
 代数 \(C\) の右noetherian性と、全てのcoherent \(\mathcal F\)、\(q\ge0\) に対する
 \[
 Q^q(\mathcal F)=\bigoplus_{\mathbf d\in\mathbb N^r}
 H^q(X,\mathcal F\mathcal L^{\mathbf d})
 \]
 の有限生成性を、\(r\) について同時に示す。
 \(r=0\) の場合は有限次元ベクトル空間の主張である。
 \(r-1\) 個の因子まで成立すると仮定し、
 \(C_{\widehat i}=\bigotimes_{j\ne i}C_j\)、\(X_{\widehat i}=\prod_{j\ne i}E_j\) と置く。
 座標 \(d_i=a\) を固定した \(Q^q(\mathcal F)\) の面は、\(C_{\widehat i}\)-加群である。
 射影 \(\pi_i:X\to X_{\widehat i}\) と
 \(\mathcal F_i=\mathcal F\mathcal L_i^a\) に対するLerayスペクトル系列の
 全多重次数の和は
 \[
 E_2^{p,b}=Q^p_{X_{\widehat i}}
       (R^b\pi_{i*}\mathcal F_i)
 \ \Longrightarrow\
 \bigoplus_{\mathbf d_{\widehat i}\ge0}
 H^{p+b}(X,\mathcal F_i\mathcal L_{\widehat i}^{\mathbf d_{\widehat i}})
 \]
 となる。射影公式と残るtranslationに関する自然性により、これは
 \(C_{\widehat i}\)-加群のスペクトル系列である。
 帰納法で各 \(E_2\) 項はnoetherianな \(C_{\widehat i}\) 上有限生成である。
 従って、有限個の核・商・拡張で得られる各面もnoetherianである。

 次の帯に関する観察を使う。
 \(\mathbb N^r\)-次数付き \(C\)-加群 \(M\) の各面 \(M_{d_i=a}\) が
 \(C_{\widehat i}\) 上noetherianであり、\(M_{\ge\mathbf b}=0\) とする。
 \[
 M^{(i)}=\bigoplus_{d_1\ge b_1,\ldots,d_i\ge b_i}M_{\mathbf d}
 \quad(0\le i\le r)
 \]
 は \(M=M^{(0)}\supset\cdots\supset M^{(r)}=0\) という \(C\)-部分加群のfiltrationである。
 各 \(M^{(i-1)}/M^{(i)}\) をさらに \(d_i\ge a\) でfiltrationすると、
 有限個の商は固定面の \(C_{\widehat i}\)-部分加群になる。
 その商では \((C_i)_{>0}\) の作用は零なので、これらは \(C\)-加群としてもnoetherianである。
 よって \(M\) はnoetherianである。この観察は \(C\) 自身のnoetherian性を使わない。
 各面がnoetherianな加群の部分商についても、同じ観察を適用できる。

 \emph{切断加群.}
 \(q>0\) のとき、amplenessにより \(Q^q(\mathcal F)_{\ge\mathbf b}=0\) となる
 \(\mathbf b\) がある。従って前段階から
 \(Q^q(\mathcal F)\) は\emph{noetherian}である。
 特に、後で用いる \(Q^1\) の任意の部分加群も有限生成である。

 \(q=0\) について、\(\mathcal F\mathcal L^{\mathbf a}\) を大域生成する
 \(\mathbf a\ge0\) を選び、逆bimoduleをtensorして
 \[
 0\longrightarrow\mathcal K\longrightarrow
       (\mathcal L^{-\mathbf a})^{\oplus t}
       \longrightarrow\mathcal F\longrightarrow0
 \]
 を得る。\(\mathcal K\) はcoherentである。
 この積の場合、\(\mathbf v\) の一つでも成分が負なら
 \(H^0(X,\mathcal L^{\mathbf v})=0\) である。
 実際、外部積のそのfactorの次数が負であり、K\"unneth公式を適用できる。
 従って切断を全非負多重次数で取った列は
 \[
 C(-\mathbf a)^{\oplus t}\longrightarrow Q^0(\mathcal F)
       \longrightarrow Q^1(\mathcal K)
 \]
 と正確である。最初の写像の像は有限生成であり、余核は既にnoetherianと示した
 \(Q^1(\mathcal K)\) の部分加群なので、\(Q^0(\mathcal F)\) は有限生成である。
 また任意の \(\mathbf a\ge0\) について
 \[
 Q^0(\mathcal F)_{\ge\mathbf a}(\mathbf a)
 \cong Q^0(\mathcal F\mathcal L^{\mathbf a})
 \]
 であるから、第一象限によるtruncationも有限生成である。

 \emph{イデアルの有限生成.}
 多重斉次右イデアル \(I\subset C\) を取る。
 \cite[Theorem 3.4]{ChanMulti} により、sheafificationはexactであり、
 \(\widetilde I\subset\widetilde C=\OO_X\) となる。
 \(X\) はnoetherianなので、このquasi-coherent部分層はcoherentである。
 随伴のunit \(\eta_I:I\to Q^0(\widetilde I)\) の自然性と
 \(\eta_C:C\xrightarrow{\sim}Q^0(\OO_X)=C\) により、\(\eta_I\) は単射である。
 その余核は多重tailsに関するtorsionであり、\(Q^0(\widetilde I)\) の商だから有限生成である。
 有限生成torsion加群は一つの第一象限上で零となるため、ある \(\mathbf a\ge0\) に対して
 \[
 I_{\ge\mathbf a}=Q^0(\widetilde I)_{\ge\mathbf a}
 \]
 となる。右辺の有限生成性は前段階で証明した。
 一方、\(I/I_{\ge\mathbf a}\) は \(C/C_{\ge\mathbf a}\) の部分加群である。
 \(C\) の固定面は有限次元空間 \((C_i)_a\) と \(C_{\widehat i}\) のtensor productなので
 \(C_{\widehat i}\)-noetherianである。
 従って帯の観察により \(C/C_{\ge\mathbf a}\) はnoetherianであり、
 \(I/I_{\ge\mathbf a}\) も有限生成となる。これで \(I\) の有限生成性を得た。

 加法と両立する \(\mathbb N^r\) の整列順序、例えば全次数付き辞書式順序で
 任意の右イデアルの先頭多重斉次項を取る。
 その斉次イデアルの有限生成元を元のイデアルへ持ち上げ、
 先頭次数を下げる帰納法を行えば、任意の右イデアルも有限生成である。
 これで通常の右noetherian性と帰納法が完了する。
 opposite algebraには逆translationを伴う正次数の可逆bimoduleが対応するので、
 同じ証明により左noetherian性も得る。
\end{proof}

\begin{lemma}[正次数中心正則元からのnoetherian性の持ち上げ]
 \label{lem:central-noetherian-lift}
 \(D\) を連結局所有限な非負次数付き代数、
 \(g\in Z(D)\) を正次数の斉次正則元とする。
 \(D/(g)\) がnoetherianなら \(D\) もnoetherianである。
\end{lemma}
\begin{proof}
 \(g\)-adic associated gradedは \((D/(g))[t]\) でありnoetherianである。
 任意の斉次右イデアル \(J\) に誘導された \(g\)-adic filtrationの
 associated gradedはこの環の右イデアルなので有限生成である。
 その二重斉次生成元を \(J\) に持ち上げる。
 内部次数 \(d\) の元のinitial termをこれらの持ち上げで消すと、
 \(g\)-adic orderが真に増える。
 しかし \(g^nD\) の内部次数は少なくとも \(n\deg(g)\) なので、
 固定した \(d\) ではこの操作は有限回で終わる。
 よって全ての斉次右イデアルが有限生成である。
 任意の右イデアルについては通常の内部次数の最高次項を取り、
 生成元を持ち上げて次数に関する帰納法を行えば有限生成性を得る。
 左側も同じである。
\end{proof}

\begin{proposition}[周囲のテンソル積]
\label{prop:tensor-sklyanin-ambient}
 \(h_i^6\ne-27\)、\(S_i=S_{h_i}\)、\(g_i=G_{h_i}\) とし、
 \[
 A=(\bigotimes_{i=1}^rS_i)[u_1,\ldots,u_m],\qquad
 N=3r+m,\qquad W=\sum_i g_i+P(u),\qquad R=A/(W)
 \]
 と置く。ここで \(P\) は任意の中心斉次三次式であり、\(m=0\) なら \(P=0\) とする。
 \(A\) は左右noetherianなKoszul AS-regular代数で、
 \[
 \operatorname{gldim}A=N,\qquad H_A(t)=(1-t)^{-N},\qquad
 \operatorname{Ext}^j_A(\CC,A)=
 \begin{cases}\CC(N)&j=N,\\0&j\ne N\end{cases}
 \]
 を満たす。\(W\) は中心正則元であり、\(R\) はinjective dimension \(N-1\)、
 Gorenstein parameter \(N-3\) のAS-Gorenstein代数である。
 この主張に \(\sigma_i\) の有限位数や \(P\) の孤立性は必要ない。
\end{proposition}
\begin{proof}
 \(S_i/(g_i)\) の斉次基底の持ち上げ \(V_i\subset S_i\) を選ぶと、
 \(g_i\) の正則性と次数に関する帰納法により
 \(S_i\cong V_i\otimes\CC[g_i]\) は中心 \(\CC[g_i]\)-加群として自由である。
 \(g_1,\ldots,g_r\) は \(\bigotimes_iS_i\) 上の中心正則列をなし、その商は
 補題\ref{lem:elliptic-thcr-tensor-noeth}の \(C\) である。
 補題\ref{lem:central-noetherian-lift}により、正次数中心正則元について
 商のnoetherian性を元の連結次数付き代数へ持ち上げられる。
 これを反復し、中心変数についてのHilbert基底定理を用いると \(A\) は左右noetherianとなる。

 各 \(S_i\) の長さ三のKoszul双加群分解と、多項式環のKoszul双加群分解をtensorすると、
 \(A\) の対角双加群の長さ \(N\) の有限射影分解を得る。
 よって \(\operatorname{gldim}A\le N\) である。
 各factorの \(\CC\) のKoszul分解とそのdualのK\"unneth公式は、
 表示したExt公式と左側の同じ式を与える。
 従ってglobal dimensionはちょうど \(N\) であり、AS-Gorenstein条件も成立する。
 Koszul性とHilbert seriesはtensor分解から従い、GK dimensionは \(N\) である。

 \(A\) は \(\CC[g_1,\ldots,g_r,u_1,\ldots,u_m]\) 上自由である。
 この中心多項式環で \(W\) は \(g_1\) に関してmonic一次式なので、\(A\) 上でも左右正則である。
 中心正則元によるchange of ringsにより、商のinjective dimensionは一つ減り、
 Gorenstein parameterは三減る。
\end{proof}

\subsection{多ブロックHesse模型からの埋込み}

\(r\ge1\)、\(m\ge0\) とし、独立した変数 \(x_j,y_j,z_j\) に対して
\[
 S_{h_j}=\CC\langle x_j,y_j,z_j\rangle/
 ([x_j,y_j]-h_jz_j^2,[y_j,z_j]-h_jx_j^2,[z_j,x_j]-h_jy_j^2),
\]
\[
 G_j=x_j^3+y_j^3+z_j^3-2h_j^2\sym(x_jy_jz_j),\qquad
 F_j=x_j^3+y_j^3+z_j^3+h_j^2x_jy_jz_j
\]
と置く。異なるブロックの元は可換とし、
\[
 A_{\mathbf h,m}=(S_{h_1}\otimes\cdots\otimes S_{h_r})[u_1,\ldots,u_m],
 \qquad W_{\mathbf h,m}=\sum_{j=1}^rG_j+\sum_{a=1}^m u_a^3,
\]
\[
 F_{\mathbf h,m}=\sum_{j=1}^rF_j+\sum_{a=1}^m u_a^3,
 \qquad Y_{\mathbf h,m}=V(F_{\mathbf h,m})\subset\PP^{3r+m-1}
\]
とする。

\begin{proposition}[多ブロックの完全忠実比較]\label{prop:multiblock-embedding}
 全ての \(j\) について \(h_j^6\ne-27\) と仮定する。
 degree shiftと三次periodicityを保つdg完全忠実関手
 \[
 \MF^{\ZZ}(\CC[x_j,y_j,z_j,u_a],F_{\mathbf h,m})
 \longrightarrow
 \MF^{\ZZ}(A_{\mathbf h,m},W_{\mathbf h,m})
 \]
 が存在する。\(n=3r+m-2\ge3\) なら、左辺は
 \((\Ku(Y_{\mathbf h,m}),\rho_{Y_{\mathbf h,m}})\) である。
 この命題は、右辺の全対象が像に属するとは主張しない。
\end{proposition}
\begin{proof}
 \(s=r+m\) とし、各非可換ブロックと各一変数ブロックに別の \(e_i\) を割り当てる。
 系\ref{cor:block}は各 \(j\) について
 \[
 \MF^{\ZZ}(\CC[x_j,y_j,z_j],F_j)
 \simeq\MF^{\ZZ}(S_{h_j},G_j)
 \]
 をdegree shiftとperiodicityを保って与える。
 これらをtensorし、\(\MF^{\ZZ}(\CC[u_a],u_a^3)\) には恒等を用いる。
 非可換側へは補題\ref{lem:multiblock-ff}の完全忠実関手を用いる。

 可換側の外部積は本質的全射でもある。
 これは可換なsmooth affine ambientに対するequivariant Thom--Sebastiani定理
 \cite[Theorem 5.15, Corollary 5.18]{BFK}を適用したものである。
 二因子の場合は \(G=H=\mathbb G_m\)、potentialのcharacterを双方 \(t\mapsto t^3\)
 と取る。異なる二因子上でのcharacterの差 \((-3,3)\) は非torsionであり、
 \(G\times_{\mathbb G_m}H\) のcharacter groupは
 \(\ZZ^2/\langle3e_2-3e_1\rangle\) となる。
 各potentialは斉次なので、Euler公式により和のcritical locusは
 個々の零点の積に含まれる。従って同定理の支持条件を除くことができる。
 反復して \(\Gamma_s\) を得る。
 ここで引用している定理のambientは全て可換であり、非可換ambientへ拡張して引用してはいない。

 これで \(\Gamma_s\)-次数付きの完全忠実関手を得た。
 \(\Gamma_s\to\ZZ\)、\(e_i\mapsto1\) の核は
 \(T_s\cong(\ZZ/3)^{s-1}\) である。
 細分次数の圏から通常次数の圏へ戻す双対同変化では、
 \(T_s\) の生成元は因子のdegree shiftの差
 \(\rho_i\rho_1^{-1}\) として作用する。
 全ての比較はこれらと三次periodicityを保つので、
 有限群関係とそのcoherenceを保つ同変関手となる。
 同変対象間のHomは元のHomの有限群不変部分であり、標数零では不変部分を取る操作は完全である。
 従って完全忠実性は同変化で保たれる。
 \cite[Theorem 4.2 and Section 8]{Elagin}の双対同変化を用いて通常次数へ降ろせば、
 求める関手を得る。\((e_1)\) との両立は降下後の \((1)\) との両立になる。

 最後に \(Y_{\mathbf h,m}\) のsmoothnessを確認する。
 特異点では全ての \(u_a=0\) であり、各三変数ブロックの偏微分が同時に零となる。
 \(F_j\) の射影零点がsmoothなので、そのaffine critical locusは原点だけである。
 従って全座標が零となり、射影特異点は存在しない。
 Orlov同値により最後の主張を得る。
\end{proof}

\begin{remark}
 命題\ref{prop:tensor-sklyanin-ambient}により、この独立ブロック族の周囲は
 平行移動の位数によらずnoetherian AS-regularである。
 中心正則元の商についてのMF--特異圏比較とOrlov windowにより、
 右辺は \(A_{\mathbf h,m}/(W_{\mathbf h,m})\) の残余成分と同一視できる。
 この周囲の正則性は、残余圏そのもののpropernessや比較の本質的全射性とは別である。
\end{remark}

\subsection{Euler導分と中心支持への帰着}

\begin{lemma}[ブロックpotentialの零ホモトピー]\label{lem:block-euler}
 任意の有限自由factorization \(P\) of \(W_{\mathbf h,m}\) 上で、
 各 \(G_j:P\to P(w)\) と各 \(u_a^2:P\to P(2e_{r+a})\) は零ホモトピックである。
 この主張は \(\Gamma_s\)-次数付きでも通常の整数次数付きでも成り立つ。
\end{lemma}
\begin{proof}
 \(j\) 番目のブロックの次数を数え、他のブロックと \(u_a\) では零となるEuler導分を
 \(\delta_j\) と書く。Sklyanin関係式はそのブロックについて斉次なので、
 \(\delta_j\) は自由代数から商代数へ降りる。
 \(\delta_j(W_{\mathbf h,m})=3G_j\) である。
 factorizationの有限自由加群に斉次基底を選び、その微分行列 \(d\) の係数に
 \(\delta_j\) を作用させる。積のLeibniz則から
 \[
 d\,\delta_j(d)+\delta_j(d)\,d
 =\delta_j(d^2)
 =3G_j\,\Id_P.
 \]
 従って \(\frac13\delta_j(d)\) が所要のホモトピーである。
 \(\delta_j\) は代数の元の \(\Gamma_s\)-次数を変えない。
 \(3e_j=w\) という次数群の関係があっても、
 元のtensor factorsから定義した導分の存在は失われない。
 同じ計算を中心変数の導分 \(\partial_{u_a}\) に対して行えば
 \[
 d\,\partial_{u_a}(d)+\partial_{u_a}(d)\,d=3u_a^2\Id_P
 \]
 を得る。基底に依存するのはホモトピーの選択であり、零ホモトピックという結論ではない。
\end{proof}

\begin{proposition}[生成性の中心商への帰着]\label{prop:block-support}
 全ての \(h_j^6\ne-27\) とし、\(A=A_{\mathbf h,m}\) と置く。
 命題\ref{prop:tensor-sklyanin-ambient}により、この \(A\) はnoetherian AS-regularである。さらに
 \[
 R=A/(W_{\mathbf h,m}),\qquad
 C_j=S_{h_j}/(G_j),\qquad C=\bigotimes_{j=1}^r C_j
 \]
 と置く。\(C\) は \(R\) の商とみなす。
 次の等式が、冪等完備化した \(\Gamma_s\)-次数付き特異圏で成り立つ。
 \[
 \Dsg^{\Gamma_s}(R)
 =\thick\!\left(
 \operatorname{Res}\bigl(D^b(\operatorname{gr}^{\Gamma_s}C)\bigr)
 \right).
 \tag{\(\ddagger\)}\label{eq:block-support}
 \]
 また、\(C\) は \(R\)-perfectなので、restrictionは特異圏の間の関手にも降りる。
\end{proposition}
\begin{proof}
 各 \(G_j\) は中心正則元である。\(S_{h_j}/G_jS_{h_j}\) の斉次基底を持ち上げた
 ベクトル部分空間を \(V_j\subset S_{h_j}\) とすると、次数に関する帰納法と正則性から
 \[
 V_j\otimes_\CC\CC[t_j]\xrightarrow{\sim}S_{h_j},\qquad
 v\otimes t_j^b\longmapsto vG_j^b
 \]
 は中心 \(\CC[t_j]\)-加群の同型である。
 従って \(A\) は中心多項式環
 \(\CC[t_1,\ldots,t_r,u_1,\ldots,u_m]\) 上自由である。
 \(W\) はこの環の \(\sum_jt_j+\sum_au_a^3\) に対応するので正則であり、
 \(R\) は \(\CC[t_1,\ldots,t_{r-1},u_1,\ldots,u_m]\) 上自由となる。
 よって
 \[
 u_1,\ldots,u_m,G_1,\ldots,G_{r-1}
 \]
 は \(R\) の中心正則列で、その商は \(C\) である。
 Koszul分解は \(C\) の有限自由 \(R\)-分解を与える。
 これでrestrictionがperfect対象をperfect対象へ送ることも従う。

 まず冪等完備化前の \(\Dsg^{\Gamma_s}(R)\) の対象 \(N\) を取る。
 中心正則元によるAS-regular代数の商に対するMF比較
 \cite[Theorem 5.8]{CCKM}を使えば、\(N\) はfinite-free factorizationのcokernelのshiftで表せる。
 この比較は通常の整数次数で記述されているが、
 \(\Gamma_s\simeq\ZZ\oplus\Gamma_{s,\mathrm{tors}}\) と標数零を用い、
 各有限自由分解とホモトピーを斉次成分ごとに取る同じ構成で
 \(\Gamma_s\)-次数を保った比較を得る。
 補題\ref{lem:block-euler}より、\(N\) は
 \[
 K(G_1,\ldots,G_{r-1},u_1^2,\ldots,u_m^2;N)
 \]
 のshiftと内部twistの直和因子となる。
 ここで \(K\) は通常の中心Koszul複体である。
 この複体の各cohomologyは \(G_1,\ldots,G_{r-1},u_1^2,\ldots,u_m^2\) に殺される。
 さらに \((u_1,\ldots,u_m)^{m+1}\subset(u_1^2,\ldots,u_m^2)\) なので、
 各cohomologyの有限filtrationの商は全 \(u_a\) と全 \(G_j\) に殺される \(C\)-加群である。
 \(G_r\) の消滅は \(W=0\) から従う。
 noetherian性によりこれらは有限生成であり、有限個のtruncation triangleとfiltrationから
 未完備の特異圏の各対象が右辺に入る。
 最後に直和因子閉包を取れば、\eqref{eq:block-support} を得る。
\end{proof}

\subsection{必要な生成性の正確な内容}

命題\ref{prop:block-support}は、全ての \(C\)-加群のrestrictionが
\(C_1,\ldots,C_r\) の外部積から生成されることまでは述べていない。
十分条件の一つは
\[
 D^b(\operatorname{gr}^{\Gamma_s}C)
 =\thick\left\{
 \operatorname{coarsen}(M_1\boxtimes\cdots\boxtimes M_r)(\gamma)
 :M_j\in D^b(\operatorname{gr}^{\ZZ}C_j),\ \gamma\in\Gamma_s
 \right\}
 \tag{G}\label{eq:block-generation}
\]
である。\(\gamma\) の全shiftは、中心変数が零となった後にも残る次数類を含めるために必要である。
以下で用いるsyzygyへの置換を正当化するため、
\(A_{\ne i}=(\bigotimes_{j\ne i}S_{h_j})[u_1,\ldots,u_m]\) とし、
他因子の外部積に全ての \(u_a\) が零で作用する \(A_{\ne i}\)-加群を \(M_{\ne i}\) とする。
\(S_{h_i}\) の \(\CC[G_i]\) 上の自由性と \(W=0\) による \(G_i\) の消去から、
\(R\) は \(A_{\ne i}\) 上自由であり、
\[
 \operatorname{Res}(C_i\boxtimes M_{\ne i})
 \simeq M_{\ne i}\otL_{A_{\ne i}}R.
\]
各 \(S_{h_j}\) 上の有限自由分解と \(u_a\) のKoszul分解を外部tensorすると、
\(M_{\ne i}\) は \(A_{\ne i}\)-perfectであり、上のrestrictionは \(R\)-perfectである。
従って一因子が \(C_i\)-perfectである項はrestriction後に零となり、
各因子を有限syzygyへ置き換えてよい。有界複体の場合もtruncation triangleから従う。
各因子の有限syzygyを取りMFで表した後、外部tensor productの有限ambient分解を
\(W\) の零ホモトピーで周期化すると、そのrestrictionの安定化は外部積factorizationになる
（全体のcohomological shiftは生成性に影響しない）。
したがって (G) が成立すれば、命題\ref{prop:block-support}と
補題\ref{lem:multiblock-ff}により非可換側の外部積もMorita同値となり、
命題\ref{prop:multiblock-embedding}を圏同値へ強められる。
より弱く、(G)の等式をrestriction後の \(\Dsg^{\Gamma_s}(R)\) だけで示しても十分である。

しかし (G) を、multigraded tailsの可換な積の幾何だけから推論することはできない。
\(C\) は元来 \(\ZZ^r\)-次数を持つが、ここで必要なのは
\(\Gamma_s\)-次数付き加群である。\(\ZZ^r\to\Gamma_s\) の核は
\(3(e_i-e_1)\) からなる無限格子を含み、有限群降下ではない。
\(\Gamma_s\)-次数付き加群は一般には \(\ZZ^r\)-次数へ持ち上がらない。
最も単純には、\(\CC[a,b]\) に \(\deg a=e_1\)、\(\deg b=e_2\)、
\(3e_1=3e_2\) を入れると、\(\CC[a,b]/(a^3-b^3)\) は \(\Gamma_2\)-次数付きだが
\(\ZZ^2\)-次数へ持ち上がらない。そのannihilatorが \(\ZZ^2\)-斉次でないためである。
この例は (G) 自体を反証するものではなく、個々の対象を細分次数へ持ち上げる証明法を反証する。

さらにmultigraded tailsでは各座標軸上のtorsionが消える。
生成性を完成させるには、これらのtorsionと上記の非liftable加群を、
restriction後も含めて制御する必要がある。
第\ref{sec:finite-center}節では、高々一つを除く因子の有限中心性でこの両方を扱い、
一つの任意位数ブロックを許して本質的全射性を証明する。
二つ以上の無限位数因子の一般の場合は未解決として保持する。

\begin{remark}[別の十分条件]
 \(\mathcal A\to\cD\) をsmooth proper圏からのdg完全忠実関手とする。
 \(\cD\) がproperでSerre関手を持ち、埋込みがSerre関手と両立すれば、
 \(\mathcal A\) のsaturated性から左右随伴が存在し、その直交成分は左右一致する。
 したがって \(\cD\) に非自明な直交直和分解がなければ埋込みは同値である。
 これは上の生成性とは別の完成条件となる。
 ただし複数ブロックの全MFについて、必要なpropernessと直交分解不能性を
 ここで証明したわけではない。
\end{remark}

\section{一方有限中心の外部積生成性と多ブロックの可換モデル}\label{sec:finite-center}

この節では生成条件 \eqref{eq:block-generation} を、高々一つを除くブロックが
有限位数である場合に証明する。任意位数の一因子は係数代数として残し、
有限中心の因子を一つずつ追加する。
非可換係数付きBeilinson分解、中心のNoether正規化、分離的な稠密開集合と
その補集合の帰納法を用いる。
\(\ZZ^r\)-次数へ持ち上がらない対象と軸上のtorsionも初めから含む。
有限中心の場合の基礎的な局所化・生成法については \cite{ELS} と比較されたい。

\subsection{斉次中心局所化と稠密開集合}

\(\operatorname{gr}^{\Lambda}D\) は有限生成の \(\Lambda\)-次数付き右 \(D\)-加群の圏とする。
ここで \(\Lambda\) は有限生成アーベル群であり、torsionを持っていてよい。
全てのthick閉包は、有限直和、shift、cone、直和因子に関する閉包を表す。

\begin{lemma}[斉次中心局所化]\label{lem:graded-localization}
 \(D\) をnoetherianな \(\Lambda\)-次数付き代数、
 \(f\in Z(D)\) を斉次元とする。
 有界導来圏の局所化は
 \[
 D^b(\operatorname{gr}^{\Lambda}D)\big/
 D^b_{f\text{-tor}}(\operatorname{gr}^{\Lambda}D)
 \simeq D^b(\operatorname{gr}^{\Lambda}D_f)
 \]
 を与える。左辺の核は、各cohomologyが \(f\) の冪に殺される対象からなる。
 この核は \(D/(f)\) 上の有限生成 \(\Lambda\)-次数付き加群のrestrictionで生成される。
\end{lemma}
\begin{proof}
 局所化は中心元についての平坦な局所化であるからcohomologyと可換である。
 有限生成加群が局所化後に零なら、その有限個の生成元を同時に殺す \(f\) の冪がある。
 \(f\)-adic filtrationと有界truncation triangleにより核についての主張を得る。

 局所化後の有限生成次数付き加群は、有限個の斉次生成元の
 \(D\)-spanを取ることによって、局所化前の有限生成次数付き加群から得られる。
 射については有限生成加群の有限生成自由分解の必要な有限部分を使う。
 中心局所化と \(\operatorname{Ext}^q\) は可換であり、
 次数零の射は \(f\) の冪を分母としてclearできる。
 同じことは有界複体にも有限個のtruncationによって成立する。
 次数を保つ分母の射は \(f^N:M(-N\deg f)\to M\) と書く。
 そのconeは \(f\)-torsionである。
 従って上記Verdier商は局所化後の圏へ完全忠実に写り、
 加群とその有限個のconeを全て含むので本質的全射である。
 この議論で次数を細分する必要はない。
\end{proof}

\begin{lemma}[有限中心代数の良い稠密開集合]\label{lem:generic-separable}
 \(k\) を標数零の体とする。
 \(D\) を連結非負次数付き \(k\)-代数で、有限生成の可換次数付き中心部分代数
 \(Z\subset Z(D)\) 上有限なものとする。
 \(Z\) がreducedであると仮定する。
 ある斉次元 \(f\in Z\) と斉次両側イデアル \(I\subset D\) が存在し、次を満たす。
 \begin{enumerate}
 \item \(f\) は \(Z\) の全てのminimal primeを避ける。
       特に \(Z/(f)\ne0\) なら \(\dim Z/(f)<\dim Z\) である。
 \item \(Z_f\) は \(k\) 上smoothな有限型可換環である。
 \item \(I_f\) は冪零であり、\(Q_f=(D/I)_f\) は
       \(Z_f\) 上有限射影かつ分離的な代数である。
 \end{enumerate}
 \(\dim Z=0\) の場合には \(Z=k\)、\(f=1\) と取れる。
\end{lemma}
\begin{proof}
 \(K\) を \(Z\) の全商環とする。
 \(Z\) はreduced noetherianなので、\(K\) はminimal primeの分数体の有限積である。
 \(D_K\) は \(K\) 上有限次元であり、そのJacobson radical \(J\) は冪零である。
 \(I\) を \(D\to D_K\) による \(J\) の逆像とする。
 次数を与える \(\mathbb G_m\)-作用はminimal primeの集合と \(J\) を保つので、
 \(I\) は斉次である。
 \(I^N\) は十分大きい \(N\) で \(K\) 上零となり、
 有限生成性から、ある全商環の分母を反転すれば \(I^N=0\) となる。

 \(Q_K=D_K/J\) は半単純である。
 標数零なので、各simple factorの中心は対応する成分体の有限分離拡大であり、
 \(Q_K\) は \(K\) 上分離的である。
 generic freenessを用いて、ある稠密開集合上で \(Q\) を有限射影とする。
 分離性は、積の写像 \(Q\otimes_ZQ^{\mathrm{op}}\to Q\) の
 \(Q\)-双加群としてのsection、すなわち分離冪等元の存在で表せる。
 generic fiberにある分離冪等元と、それが満たす有限個の等式の分母を払えば、
 さらに小さい稠密開集合上でこの分離冪等元が存在する。
 また \(Z\) のsmooth locusは全てのgeneric pointを含む。

 以上の良い開集合は \(\mathbb G_m\)-作用で不変に取れる。
 実際、\(I^N\) の消滅、局所自由性、分離性、smoothnessは全て
 次数作用によって保たれる開条件である。
 補集合のイデアルは斉次で、どのminimal primeにも含まれない。
 \(Z_0=k\) と \(k\) の無限性により、そのイデアル内から
 全てのminimal primeを避ける一つの斉次元 \(f\) を選べる。
 必要なら候補を冪にして次数を揃え、一般の線型結合を取ればよい。
 \(\dim Z>0\) のとき零次元の既約成分はない
 （全ての既約成分は連結な錐の頂点を含む）ので、この操作に例外はない。
 \(Z=k\) の場合は有限次元代数のradicalと半単純商を用いればよい。
\end{proof}

\begin{lemma}[分離代数と射影分解]\label{lem:separable-perfect}
 \(Z\) をsmooth affineな有限型 \(k\)-代数、
 \(Q\) を \(Z\) 上有限射影かつ分離的な代数とする。
 このとき \(Q\) は有限global dimensionを持つ。
 特に \(Q\) に任意のアーベル群次数を与えたとき、
 \(D^b(\operatorname{gr}^{\Lambda}Q)\) は \(Q(\lambda)\)
 \((\lambda\in\Lambda)\) で生成される。
\end{lemma}
\begin{proof}
 任意の右 \(Q\)-加群 \(M\) は、分離冪等元により
 \(M\otimes_ZQ\to M\) の \(Q\)-線型sectionを持つ。
 \(Z\) は有限global dimensionを持ち、\(Q\) は \(Z\)-flatなので、
 \(M\) の \(Z\)-射影分解を \(-\otimes_ZQ\) で送ると
 \(M\otimes_ZQ\) の有限 \(Q\)-射影分解を得る。
 従ってその直和因子 \(M\) も有限射影次元を持つ。

 次数付き有限生成加群に有限生成斉次自由分解を作ると、十分高いsyzygyは
 ungradedには射影的である。
 自由加群からこのsyzygyへの次数零の全射のungraded sectionを取り、
 sectionの次数零成分を抜き出すと次数付きsectionになる。
 よってsyzygyは次数付き有限自由加群の直和因子であり、
 元の加群は有限個の \(Q(\lambda)\) から生成される。
\end{proof}

\subsection{非可換係数の多項式拡大と一方有限中心の生成定理}

ここでは係数代数の有限中心性を仮定しない。
これが一つの無限位数ブロックを許すための要点である。
以下、\(k\) は標数零の体、\(\Lambda\) は有限生成アーベル群とする。

\begin{lemma}[非可換係数を持つ多項式拡大]\label{lem:polynomial-coefficient}
 \(A\) を左右noetherianな \(\Lambda\)-次数付き \(k\)-代数とする。
 \(P=k[t_1,\ldots,t_d]\) の各変数に任意の \(\Lambda\)-次数を与え、
 \(A\otimes_kP\) には和の次数を与える。このとき
 \[
 D^b(\operatorname{gr}^{\Lambda}(A\otimes_kP))
 =\thick\{M\otimes_kP:\ M\in\operatorname{gr}^{\Lambda}A\}.
 \]
 \(P\) の斉次元 \(f\) を反転した \(P_f\) についても同じ結論が成立する。
 \(A\) が有限global dimensionを持つという仮定はない。
\end{lemma}
\begin{proof}
 中心多項式拡大のHilbert基底定理により \(A\otimes P\) は左右noetherianである。
 \(d=0\) は自明なので \(d>0\) とする。
 有限生成 \(\Lambda\)-次数付き \(A\otimes P\)-加群 \(T\) の有限斉次表示を取る。
 射影用の独立な整数次数を導入し、\(t_0,t_1,\ldots,t_d\) の射影次数を全て1、
 \(A\) の射影次数を0とする。\(\Lambda\)-次数は
 \(\deg_\Lambda(t_0)=0\)、\(\deg_\Lambda(t_i)=\deg_\Lambda(t_i\in P)\) とする。
 表示の各関係を \(t_0\) で斉次化すると、\(\PP^d_k\) 上の
 \(\mathcal A=A\otimes_k\OO_{\PP^d}\)-加群 \(\cF\) が得られる。
 これは \(\mathcal A\) 上有限表示で、\(t_0\ne0\) への制限が \(T\) である。
 各関係の斉次化では \(\Lambda\)-次数を変えていない。
 以後の \(\Lambda\)-次数は対角化可能群 \(D(\Lambda)\) による線型化としても解釈できる。

 必要なcohomologyの有限性を確認する。
 \(\mathcal A\) 上のcoherent加群には有限個の
 \(\mathcal A(-a)(\lambda)\) からの全射があり、核もcoherentである。
 これは中心変数による有限斉次表示とnoetherian性から従う。
 標準affine被覆の\v Cech複体は長さ \(d\) なので、cohomologyは次数 \(0,\ldots,d\)
 に限られる。一方、\(H^q(\PP^d,\mathcal A(a))
 =A\otimes_kH^q(\PP^d,\OO(a))\) は有限生成 \(A\)-加群である。
 上記全射とその核の長完全列を用い、\(q=d,d-1,\ldots,0\) の順に帰納すれば、
 全てのcoherent \(\mathcal A\)-加群の \(H^q\) が有限生成 \(A\)-加群と分かる。
 この議論は線型化と両立するので、\(\cF\) に有限階数線束をテンソルしたものの
 導来大域切断は \(D^b(\operatorname{gr}^{\Lambda}A)\) に属する。
 ここでcoherentは \(\mathcal A\) 上の意味であり、\(\OO\) 上有限とは仮定しない。

 Beilinsonの対角分解 \cite{Beilinson} を、因子の順序を適切に選んで
 \[
 0\longrightarrow\Omega^d(d)\boxtimes\OO(-d)
 \longrightarrow\cdots\longrightarrow
 \Omega^1(1)\boxtimes\OO(-1)
 \longrightarrow\OO\boxtimes\OO
 \longrightarrow\OO_\Delta\longrightarrow0
 \]
 と書く。これは \(D(\Lambda)\)-線型化を保つ分解である。
 \(\mathcal A\) の係数を保ったまま \(\cF\) に作用させると、\(\cF\) は
 \[
 R\Gamma(\PP^d,\cF\otimes_{\OO}\Omega^i(i))\otimes_k\OO(-i),
 \qquad 0\le i\le d,
 \]
 の有限個のconeから作られる。
 対角のKoszul分解は各射影に関してflatであるから、\(\cF\) の
 \(\OO\)-flatnessは必要ない。非可換な \(A\)-作用も全ての操作と可換である。
 係数複体は前段落により \(D^b(\operatorname{gr}^{\Lambda}A)\) に属する。
 \(t_0\ne0\) に制限すると、\(t_0\) の \(\Lambda\)-次数が0であるため
 \(\OO(-i)\) は次数を保って自明化される。
 よって \(T\) は \(M\otimes P\) から生成される。
 係数複体を有限個のcohomology加群に分けるには通常のtruncationを用いればよい。
 有界複体全体についても同じtruncationから結論する。

 \(P_f\) の場合は補題\ref{lem:graded-localization}による局所化を適用する。
 有限生成次数付き加群の有限表示は局所化前に持ち上がるので、
 局所化後の生成対象はちょうど \(M\otimes P_f\) である。
\end{proof}

\begin{lemma}[多項式中心上の分離的な稠密開集合]\label{lem:polynomial-separable}
 \(D\) を連結局所有限な非負次数付き \(k\)-代数で、有限型可換斉次中心部分代数
 \(Z\) 上有限なものとする。\(d=\dim Z\) と置く。
 正次数の斉次変数を持つ多項式部分代数 \(P=k[t_1,\ldots,t_d]\subset Z\)、
 非零斉次元 \(f\in P\)、斉次両側イデアル \(J\subset D\) を選んで、
 \(D\) が \(P\) 上有限、\(J_f\) が冪零、\(Q=(D/J)_f\) が
 \(P_f\) 上有限射影かつ分離的となるようにできる。
 分離冪等元は元の整数次数で次数0に取れる。
\end{lemma}
\begin{proof}
 斉次Noether正規化によって \(Z\) は正次数変数の多項式部分代数 \(P\) 上有限となる。
 \(Z\) のreducednessは必要ない。\(K=\operatorname{Frac}P\) とすると、
 \(D_K\) は有限次元 \(K\)-代数である。
 そのradicalの \(D\) への逆像を \(J\) とする。
 次数のスカラー作用は \(P\)、\(K\) とradicalを保つので \(J\) は斉次である。
 十分大きい \(a\) で \((J^a)_K=0\) だから、有限生成性によって
 ある非零元を反転すると \(J^a=0\) となる。
 半単純商 \(D_K/\operatorname{rad}(D_K)\) は標数零の体上分離的である。
 generic freenessと分離冪等元の有限個の等式の分母払いにより、
 ある稠密開集合上で \(D/J\) は有限射影かつ分離的となる。

 これらの条件の成立する開集合は次数作用に関して不変である。
 その補集合のイデアルは斉次であり、\(P\) は整域なので、そこから
 非零斉次元 \(f\) を一つ選べる。
 \(d=0\) の場合は \(P=k\)、\(f=1\) とし、有限次元代数のradicalを使う。
 最後に、分離冪等元を元の整数次数で分解し、次数0の成分を取る。
 \(\mu(e)=1\) と \((q\otimes1)e=e(1\otimes q)\) は斉次 \(q\) ごとに次数に分解されるから、
 次数0成分も同じ条件を満たす。これが必要な次数付き分離冪等元である。
\end{proof}

\begin{theorem}[一方だけが有限中心である外部積生成]\label{thm:one-sided-boxes}
 \(A\) を左右noetherianな \(\Lambda\)-次数付き \(k\)-代数とする。
 \(D\) を連結局所有限な非負次数付き \(k\)-代数で、
 有限型可換斉次中心部分代数 \(Z\) 上有限なものとする。
 任意の準同型 \(\beta:\ZZ\to\Lambda\) により \(D\) の次数を粗化し、
 \(A\otimes_kD\) には和の次数を与える。このとき
 \begin{equation}\label{eq:one-sided-boxes}
 \begin{split}
 D^b(\operatorname{gr}^{\Lambda}(A\otimes_kD))
 =\thick\{\,M\boxtimes\operatorname{coarsen}_\beta N:\qquad\qquad\qquad\quad\\
 M\in\operatorname{gr}^{\Lambda}A,\quad
 N\in\operatorname{gr}^{\ZZ}D\,\}.
 \end{split}
 \end{equation}
 特に \(A\) の有限中心性、有限global dimension、あるいは
 左辺の対象の細分次数へのliftは仮定しない。
\end{theorem}
\begin{proof}
 補題\ref{lem:polynomial-separable}の \(P\subset Z\) に対して
 \(A\otimes D\) は \(A\otimes P\) 上左右有限である。
 後者はnoetherianなので前者もnoetherianである。
 右辺のthick部分圏を \(\mathcal B\) とし、\(d=\dim Z\) に関して帰納する。
 \(d=0\) では \(D\) は有限次元であり、連結性から \(D_{>0}\) は冪零で商は \(k\) である。
 \(D_{>0}\)-adic filtrationの各層は有限生成 \(\Lambda\)-次数付き \(A\)-加群であるから、
 \(M\boxtimes k\) からの生成を得る。

 \(d>0\) とし、同補題の \(f,J,Q\) を取る。
 \(f=1\) と取れる場合も以下の開集合の議論はそのまま適用でき、閉集合側は空になる。
 \(J_f\)-adic filtrationにより、\(A\otimes D_f\) 上の全加群は
 \(A\otimes Q\) 上の有限生成加群から生成される。
 そのような加群 \(T\) は \(A\otimes P_f\) 上も有限生成である。
 次数0の分離冪等元 \(e=\sum_\nu x_\nu\otimes y_\nu\) により、乗法写像
 \[
 \operatorname{Res}T\otimes_{P_f}Q\longrightarrow T
 \]
 は \(A\otimes Q\)-線型かつ \(\Lambda\)-次数0のsection
 \(t\mapsto\sum_\nu tx_\nu\otimes y_\nu\) を持つ。
 \(A\) の作用は \(Q\) と可換なので、このsectionは \(A\)-線型でもある。

 補題\ref{lem:polynomial-coefficient}により \(\operatorname{Res}T\) は
 \(M\otimes P_f\) から生成される。
 \(Q\) は \(P_f\) 上射影なので \(-\otimes_{P_f}Q\) はexactであり、
 上のsectionと合わせて \(T\) は \(M\otimes Q\) から生成される。
 \(M\otimes Q\) は \(M\boxtimes(D/J)\) の局所化であるから、
 開集合上の全圏は \(\mathcal B\) の像で生成される。
 ここで \(A\otimes Q\) が有限global dimensionを持つとは主張していない。

 局所化の核は補題\ref{lem:graded-localization}により、
 \(A\otimes(D/(f))\) 上の有限生成次数付き加群から生成される。
 \(D/(f)\) は \(P/(f)\) の像という有限型可換斉次中心上有限である。
 \(f\) が非単元ならこの中心の次元は \(d-1\) 以下であるから、帰納法を適用できる。
 得られる外部積の第二因子は \(D/(f)\)-加群であり、元の \(D\)-加群としても
 元の整数次数を保っている。従って局所化の核は \(\mathcal B\) に含まれる。
 開集合と核の生成をVerdier局所化によって合わせれば \eqref{eq:one-sided-boxes} を得る。
 全てのconeと直和因子は粗化後の圏で取っており、対象の次数のliftは用いない。
\end{proof}

\subsection{任意の粗化次数と有限中心因子の反復}

\begin{theorem}[有限中心代数の外部積生成]\label{thm:finite-center-boxes}
 \(D_1,\ldots,D_r\) を連結局所有限な非負次数付き \(k\)-代数で、
 各 \(D_i\) は有限型可換斉次中心部分代数上有限とする。
 任意の準同型 \(\alpha:\ZZ^r\to\Lambda\) に対し、
 \(D=\bigotimes_iD_i\) に粗化次数を与えると
 \[
 D^b(\operatorname{gr}^{\Lambda}D)
 =\thick\{\operatorname{coarsen}_{\alpha}
 (M_1\boxtimes\cdots\boxtimes M_r)(\lambda):
 M_i\in\operatorname{gr}^{\ZZ}D_i,\ \lambda\in\Lambda\}.
 \tag{\(G_{\mathrm{fin}}\)}\label{eq:finite-center-generation}
 \]
 \(\alpha\) は単射・全射である必要はなく、無限の核を持っていてもよい。
\end{theorem}
\begin{proof}
 \(A=k\) から出発し、\(\beta_i(1)=\alpha(e_i)\) として
 定理\ref{thm:one-sided-boxes}を順に適用する。
 各段階の係数代数は有限型可換中心上有限なので左右noetherianである。
 \(k\) の有限生成 \(\Lambda\)-次数付き加群は有限個の \(k(\lambda)\) の和であり、
 反復によって表示の外部積を得る。
 複体の外部積は体 \(k\) 上exactなので、前段階のthick生成を次段階へ送れる。
\end{proof}

\begin{remark}[nonliftable対象と軸torsion]
 外部積のconeには粗化後の全ての次数零の射を用いる。
 例えば \(k[a,b]/(a^3-b^3)\) は \(3e_1=3e_2\) の次数では
 \(a^3-b^3\) による二つのfree shiftの間のconeで生成される。
 同じ加群が \(\ZZ^2\)-次数へ持ち上がる必要はない。
 また \(f\)-adic filtrationは正次数の中心元に殺される軸上の加群も保持する。
 multigraded tailsへの移行によってこれらを捨てる操作は用いない。
\end{remark}

\subsection{高々一つの無限位数ブロックへの帰結}

次の有限中心性は \cite[Corollary 2.3]{SmithTate} に含まれる既知の結果である。
必要な斉次中心を明示するため、証明を記す。

\begin{lemma}[有限位数THCRの有限中心性]\label{lem:finite-order-thcr}
 \(E\) をsmooth elliptic curve、\(L\) を正次数の線束とし、
 \(\sigma\in\operatorname{Aut}(E)\) が有限位数 \(q\) を持つとする。
 このとき \(B=B(E,L,\sigma)\) は有限型可換斉次中心部分代数上有限である。
\end{lemma}
\begin{proof}
 \(L=\mathcal O_E(D)\) となるdivisor \(D\) を選び、
 \(D_d=D+\sigma^*D+\cdots+(\sigma^{d-1})^*D\) と置く。
 通常のskew Laurent表示により、\(B\) を
 \[
 B_d=H^0(E,\mathcal O_E(D_d))t^d
 \subset k(E)[t,t^{-1};\sigma^*]
 \]
 と書ける。\(\sigma^q=1\) なので \(D_q\) は \(\sigma\)-不変であり、
 \(q\)-Veroneseは可換なsection ring
 \[
 R=B^{(q)}=\bigoplus_{n\ge0}H^0(E,\mathcal O_E(nD_q))t^{qn}
 \]
 になる。\(R\) 上では \(\sigma\) を係数に作用させ、\(t^q\) を固定する。
 \(Z=R^{\langle\sigma\rangle}\) は \(B\) の中心に含まれる。
 実際、\(a\in H^0(E,\mathcal O_E(nD_q))\) が \(\sigma\)-不変なら、
 任意の \(bt^d\in B_d\) に対して
 \[
 (at^{qn})(bt^d)=ab\,t^{qn+d}
 =b\sigma^{d*}(a)t^{qn+d}=(bt^d)(at^{qn})
 \]
 である。

 \(D_q\) は正次数なのでampleである。従って \(R\) は有限型可換代数であり、
 各 \(0\le d<q\) について
 \[
 \bigoplus_{n\ge0}H^0(E,\mathcal O_E(nD_q+D_d))t^{qn+d}
 \]
 は \(R\) 上有限なsection moduleである。
 これらの有限直和が \(B\) なので \(B\) は \(R\) 上有限である。
 有限群の不変式環について \(R\) は有限型環 \(Z\) 上有限であり、
 従って \(B\) も \(Z\) 上有限である。
\end{proof}




\begin{corollary}[一つの任意位数因子を含む生成条件]\label{cor:PI-generation}
 全ての \(j\) で \(h_j^6\ne-27\) とし、並べ替えた後の
 \(\sigma_j=t_{p_j}\ (j\ge2)\) のみ有限位数と仮定する。
 このとき \(C_j=S_{h_j}/(G_j)\)、\(C=\bigotimes_jC_j\) に対して、
 条件 \eqref{eq:block-generation} は成立する。
 中心変数を零とした後に残る \(\Gamma_s\) の全shiftも含む。
\end{corollary}
\begin{proof}
 \(C_1\) はnoetherianであり、その \(\Gamma_s\)-次数は
 単射 \(\ZZ e_1\hookrightarrow\Gamma_s\) によって与えられる。
 任意の有限生成 \(\Gamma_s\)-次数付き \(C_1\)-加群は
 \(\Gamma_s/\ZZ e_1\) の剰余類ごとの部分加群に分解する。
 有限個の生成元の次数が関わる剰余類は有限個だけであるから、
 各剰余類の代表によるshiftを戻せば、元の整数次数を持つ \(C_1\)-加群の有限和となる。
 ここで細分しているのは一因子上の加群だけであり、積の一般の加群ではない。

 \(j\ge2\) では補題\ref{lem:finite-order-thcr}によって \(C_j\) は有限中心上有限である。
 定理\ref{thm:one-sided-boxes}を
 \(A=C_1\)、\(D=C_2\)、\(\beta(1)=e_2\) に適用し、次に係数代数を
 \(C_1\otimes C_2\) として反復する。
 各段階で左右noetherian性は同定理の有限中心多項式拡大によって保証される。
 外部積は体上exactなので、前段階で得たthick生成をそのまま次段階へ送れる。
 これで全ての \(C_j\) について条件 \eqref{eq:block-generation} が得られる。
\end{proof}

以下では定理\ref{thm:multiblock-model}の仮定の下で
\[
 \begin{gathered}
 N=3r+m,\qquad A=\left(\bigotimes_{j=1}^rS_{h_j}\right)[u_1,\ldots,u_m],\\
 W=\sum_jG_{h_j}+P(u),\qquad R=A/(W)=B_{\mathbf h,P}.
 \end{gathered}
\]
と置く。\(m=0\) の場合は \(P=0\) とし、中心変数の因子を省く。

\begin{proof}[定理\ref{thm:multiblock-model}の証明]

 \emph{周囲の代数と中心元の正則性.}
 命題\ref{prop:tensor-sklyanin-ambient}により、\(A\) は左右noetherianな
 Koszul AS-regular代数で、global dimensionとGorenstein parameterは \(N\) である。
 \(W\) は中心正則元である。以下の支持帰着のために、中心上の自由性も具体的に用いる。

 \(C_j\) の斉次基底の持ち上げを \(V_j\subset S_{h_j}\) とすると、
 正則性と次数に関する帰納法により
 \[
 V_j\otimes\CC[t_j]\xrightarrow{\sim}S_{h_j},\qquad
 v\otimes t_j^a\longmapsto vG_j^a
 \]
 は中心 \(\CC[t_j]\)-加群としての同型である。
 従って \(A\) は中心多項式環
 \(Z_0=\CC[t_1,\ldots,t_r,u_1,\ldots,u_m]\) 上自由である。
 \(W\) は \(\sum_jt_j+P(u)\in Z_0\) の作用なので、左右正則である。
 その商 \(R\) は
 \(\CC[t_1,\ldots,t_{r-1},u_1,\ldots,u_m]\) 上自由であり、
 \[
 u_1,\ldots,u_m,G_1,\ldots,G_{r-1}
 \]
 は \(R\) の中心正則列で、商は \(C\) である。

 \emph{全生成性.}
 \(m>0\) では中心多項式環全体 \((T,P)=(\CC[u],P)\) を一つの因子とし、
 \[
 \Gamma=\ZZ^{r+1}/\langle3e_j-3e_0:1\le j\le r\rangle
 \]
 を用いる。\(\deg(u_a)=e_0\)、Sklyanin因子の次数は \(e_j\) である。
 \(m=0\) では \(e_0\) を省き、\(\Gamma=\Gamma_r\) とする。
 finite-free factorizationの微分を \(d\)、第 \(j\) ブロックのEuler導分を
 \(\delta_j\) とすれば、補題\ref{lem:block-euler}と同じ計算により
 \[
 d\delta_j(d)+\delta_j(d)d=3G_j\operatorname{id}.
 \]
 中心変数についても行列係数を微分することで
 \[
 d\,\partial_{u_a}(d)+\partial_{u_a}(d)\,d
 =\partial_{u_a}P\operatorname{id}
 \]
 を得る。これらは斉次のnull-homotopyである。

 \(J_P=(\partial_{u_1}P,\ldots,\partial_{u_m}P)\) とする。
 孤立性より、ある \(b\) について \((u_1,\ldots,u_m)^b\subset J_P\) である。
 命題\ref{prop:block-support}の証明で \(u_a^2\) を \(\partial_{u_a}P\) に置き換えると、
 各対象は \(G_1,\ldots,G_{r-1},\partial_{u_1}P,\ldots,\partial_{u_m}P\)
 に関するKoszul複体のshift・内部twistの直和因子になる。
 各cohomologyの \((u)\)-adic有限filtrationの商は全ての \(u_a\) に殺される。
 それらは \(G_1,\ldots,G_{r-1}\) にも殺され、\(W=0\) と \(P(0)=0\) から
 \(G_r\) にも殺される。従って
 \[
 \Dsg^\Gamma(R)
 =\thick\operatorname{Res}D^b(\operatorname{gr}^\Gamma C).
 \]
 \(m=0\) では中心変数についてのKoszul複体とfiltrationを省けばよい。
 全ての \(\Gamma\)-次数付き \(C\)-加群が
 系\ref{cor:PI-generation}により外部積で生成される。

 この外部積のrestrictionのstabilizationを説明する。
 各 \(C_j=S_{h_j}/(G_j)\) の加群を十分高いsyzygyに置き換えて二項ambient自由分解で表し、
 \(m>0\) では \(T/(P)\) 上の残余体 \(k\) についても同じ操作をする。
 一つのhypersurface因子においてperfectな項は、restriction後に \(R\)-perfectとなる。
 この点を右加群の規約で明示する。選んだambientを \(A_i\)、potentialを \(f_i\)、
 その商を \(D_i\)、他のambientの積を \(A'\)、potentialの和を \(g'\) と書く。
 \(A_i\cong V_i\otimes\CC[f_i]\) と \(f_i=-g'\) の消去により、
 \(R=(A_i\otimes A')/(f_i+g')\) は \(A'\) 上左右とも自由である。
 他因子の外部積 \(M'\) は \(M'g'=0\) を満たし、各ambient上の有限自由分解の
 tensor productによって \(A'\)-perfectである。従って
 \[
 \begin{aligned}
 M'\otL_{A'}R&\simeq M'\otimes_{A'}R
       \xrightarrow{\sim}\operatorname{Res}(D_i\boxtimes M'),\\
 m'\otimes[a\otimes b]&\longmapsto\bar a\otimes m'b
 \end{aligned}
 \]
 は次数を保つ同型であり、右辺は \(R\)-perfectである。
 有限cone、shift、直和因子を取ることで、\(D_i\) を任意の \(D_i\)-perfect複体に
 置き換えても成立する。
 中心因子についても \(m>0\) では孤立性から \(P\ne0\) であり、
 \(T\cong V_0\otimes\CC[P]\) となるので同じ議論を適用できる。
 従ってsyzygyへの置換で失う項は特異圏で零である。

 二項ambient分解の微分を \(d_i\)、potential倍のnull-homotopyを \(s_i\) とする。
 外部tensor積上のKoszul符号付きの和 \(d=\sum d_i\)、\(s=\sum s_i\) は
 \(ds+sd=W\)、\(s^2=0\) を満たす。
 そのShamash周期化のtailはfactorizationの外部積に一致する。
 よって各boxのrestrictionのstabilizationは外部積関手のthick imageに入る。
 全 \(\Gamma\)-shiftも各因子のdegree shiftで実現されるので、
 補題\ref{lem:multiblock-ff}の完全忠実関手はMorita同値である。

 \emph{可換側との比較と整数次数への降下.}
 各Sklyanin因子に系\ref{cor:block}を用い、中心因子には恒等を用いる。
 可換側の外部積には \cite[Theorem 5.15, Corollary 5.18]{BFK} を適用する。
 命題\ref{prop:multiblock-embedding}と同様に、\(\Gamma\to\ZZ\) の有限核の双対による
 同変化と双対同変化を用いる。因子ごとのdegree shiftと \((3)=[2]\) のcoherenceは
 各比較で保存されるので、通常の整数次数へ降下すると
 \[
 \MF^{\ZZ}(A,W)\simeq
 \MF^{\ZZ}\!\left(\CC[x_j,y_j,z_j,u_a],
 \sum_j(x_j^3+y_j^3+z_j^3+h_j^2x_jy_jz_j)+P(u)\right)
 \tag{\(*\)}\label{eq:PI-model}
 \]
 を得る。

 \emph{射影模型と一意性.}
 中心正則元の商 \(R\) はnoetherian AS-Gorensteinで、injective dimensionは \(N-1\)、
 Gorenstein parameterは \(N-3\) である。
 \cite[Theorem 5.8]{CCKM} のMF--特異圏比較と
 \cite[Theorem 2.5]{Orlov} の残余windowを適用できる。
 \(N>3\) ならそのwindowは
 \[
 D^b(\qgr R)=\langle\mathcal K_R,\OO,\OO(1),\ldots,\OO(N-4)\rangle
 \]
 を与え、\(N=3\) なら例外対象はなく \(\mathcal K_R=D^b(\qgr R)\) である。
 degree shiftはwindowを通してrotationに対応する。
 可換potentialの勾配は各変数群に分離しており、各 \(F_{h_j}\) と \(P\) の
 アフィン臨界点は原点だけであるから、その和も原点以外に臨界点を持たない。
 従って \(Y_{\mathbf h,P}\) は滑らかで、同じOrlov比較は可換側を
 \((\Ku(Y_{\mathbf h,P}),\rho_{Y_{\mathbf h,P}})\) と同定する。
 \(N-2\ge3\) における射影同型類の一意性は
 \cite[Theorem 1.7]{LZ} による。
\end{proof}

\begin{lemma}[位数五の明示的パラメータ]\label{lem:order-five}
 \(h^6=5\) とする。原点 \(o=[1:-1:0]\) を持つ
 \(E_h=V(x^3+y^3+z^3+h^2xyz)\) 上で、\(p=[1:-1:-h]\) はちょうど位数五である。
 従って対応するSklyaninの平行移動 \(\sigma=t_p\) も位数五である。
\end{lemma}
\begin{proof}
 \(o\) における接線は \(3x+3y-h^2z=0\) であり、
 この直線上で \(F_h=(1+h^6/27)z^3\) となるので、\(o\) はinflection pointである。
 逆元は \([x:y:z]\mapsto[y:x:z]\) である。実際、この二点と \(o\) は一直線上にある。
 \(a=h^3\) とし、\(q=[a+1:a-1:-2h]\) と置く。
 \(F_h(q)=0\) であり、\(p\) における接線
 \[
 (3+a)x+(3-a)y+2h^2z=0
 \]
 は \(-q=[a-1:a+1:-2h]\) を通る。
 \(h\ne0\) では \(-q\ne p\) であるから、接線の加法則により \(q=2p\) である。
 \(q\) における接線の係数は
 \[
 [a^2+8a+3:a^2-8a+3:h^2(a^2+11)]
 \]
 であり、これを \(p\) で評価すると \(a(5-a^2)=h^3(5-h^6)\) となる。
 \(p\ne q\) は \(q=2p\) と \(p\ne o\) から従う。
 \(h^6=5\) では \(p\) がこの接線上にあるので \(2q+p=0\)、従って \(5p=0\) である。
 \(p\ne o\) かつ五は素数なので、位数はちょうど五である。
\end{proof}

\begin{remark}[二・三ブロックのsixfold・sevenfold]
 全ての \(h_j^6=5\) とする。
 \(r=2,m=2,P=u_1^3+u_2^3\) は二つの非可換ブロックを持つcubic sixfoldの
 rotation付き残余圏の完全な可換模型を与える。
 \(r=3,m=0\) は三つの非可換ブロックを持つcubic sevenfoldの完全な可換模型を与える。
 両者の方程式は定理\ref{thm:multiblock-model}の \(Y_{\mathbf h,P}\) であり、
 元の対象の一部分ではなく全残余圏を実現する。
\end{remark}

\begin{lemma}[無限位数の明示的パラメータ]\label{lem:order-infinite}
 \(h=2\) とする。楕円曲線
 \[
 E_2=V(x^3+y^3+z^3+4xyz),\qquad o=[1:-1:0]
 \]
 上の点 \(p_2=[1:-1:-2]\) は無限位数である。
\end{lemma}
\begin{proof}
 曲線と点は \(\mathbb Q\) 上定義されている。
 \(2^6+27=91\) は5,11のどちらでも割れないので、Hesseの勾配による判定から
 5と11で良い還元を持つ。
 \(z=1\)、\(z=0,y=1\)、\([1:0:0]\) の三つのchartに分けて列挙すると
 \[
 \#E_2(\mathbb F_5)=9,\qquad \#E_2(\mathbb F_{11})=12.
 \]
 実際、\(z=1\) の点数はそれぞれ8と11、\(z=0\) にはどちらも一点だけ存在する。
 この列挙は整数の合同計算で再現でき、\texttt{verify\_algebra.py} に全座標を記録する。
 良い還元では剰余標数と素なtorsionの還元写像は単射である
 \cite[Reduction of torsion points]{SnowdenDVR}。
 従って \(p_2\) の位数 \(N\) が有限なら、\(N\) の5と素な部分は9を割り、
 11と素な部分は12を割る。後者により5-primary部分はなく、前者により11-primary部分もなく、
 残りの素数について両方を比較すれば \(N\mid3\) となる。

 \(p_2\ne o\) である。一方、接線による倍点公式から
 \[
 2p_2=[9:7:-4],\qquad -p_2=[1:-1:2]
 \]
 であり、これらは異なる射影点である。
 従って \(3p_2\ne0\) であり、\(N\mid3\) は不可能である。
 \(p_2\) は無限位数である。
\end{proof}

\begin{remark}[真に無限位数を含むsixfoldとsevenfold]
 \(h_1=2\)、\(h_j=5^{1/6}\ (j\ge2)\) と取る。
 第一ブロックの平行移動は無限位数、他のブロックは位数5である。
 \((r,m)=(2,2)\)、\(P=u_1^3+u_2^3\) では可換sixfoldは
 \[
 V\bigl(x_1^3+y_1^3+z_1^3+4x_1y_1z_1
       +x_2^3+y_2^3+z_2^3+5^{1/3}x_2y_2z_2+u_1^3+u_2^3\bigr)
 \subset\PP^7
 \]
 となる。\((r,m)=(3,0)\) では第二・第三ブロックの係数を \(5^{1/3}\) とした
 cubic sevenfoldを得る。いずれも全残余圏とrotationの同値である。
 さらに他の有限位数パラメータを固定すれば、\(h_1^6\ne-27\) を満たす
 全ての \(h_1\in\CC\) について同じ結論が成り立つ。
 これは各ファイバーでの主張であり、同値の族をこのパラメータ直線全体に
 代数的に降下させたという主張ではない。
\end{remark}

\subsection{証明範囲と残る課題}

有限位数条件は \(j\ge2\) のみに用いる。
従ってより一般に、\(C_1\) がnoetherian、\(j\ge2\) の各 \(C_j\) が
有限型可換斉次中心上有限で、周囲が命題\ref{prop:block-support}の仮定を満たせば、
同じ生成と安定化の議論を適用できる。
定理\ref{thm:one-sided-boxes}では任意のnoetherianな係数代数を保てる一方、
新たに加える因子の有限中心性は必要である。
従って二つ以上の独立な無限位数THCRを持つ場合は、この証明だけでは解けない。
これはその場合の非存在を意味せず、条件 (G) が未解決であるという意味である。

次の焦点は、二つの無限位数因子の条件 (G) を証明することである。
それが得られれば、残りの有限中心因子は定理\ref{thm:one-sided-boxes}で追加できる。
必要なのは、THCRに対しても任意のnoetherian係数を許す外部積生成、
または粗化次数を保った対角の有限構成である。
可換ambientのThom--Sebastiani定理を非可換ambientへ無条件に適用したわけではない。
先行研究上の新規性は未確認である。

\section{適用範囲と残る問題}\label{sec:scope}

\subsection{例外パラメータでpropernessが失われる}

\begin{proof}[命題\ref{prop:bad-intro}の証明]
 \(h_0^6=-27\) とする。\(E_{h_0}\) は三本の直線からなる三角形である。
 まず
 \[
 S_{h_0}/(G_{h_0})\simeq B(E_{h_0},L,\sigma_{h_0})                 \tag{6}\label{eq:triangle}
 \]
 を確認する。楕円曲線の場合だけの公式を、この特殊値へ無条件に適用しない。
 \(h_0\) の近傍で、関係式を多重線形化した行列
 \[
 M_h(x,y,z)=
 \begin{pmatrix}-y&x&-hz\\-hx&-z&y\\z&-hy&-x\end{pmatrix}
 \]
 は \(E_h\) 上でrank二である。実際その行列式は
 \(-h(x^3+y^3+z^3+h^2xyz)\) であり、rank一なら小行列式から
 \(yz+hx^2=yz-hx^2=0\) と巡回的な式を得て、\(x=y=z=0\) となる。
 核は相対的な射 \(\sigma_h\) を与える。グラフを転置すると \(h\) を \(-h\) に
 変える同じ関係式になるので、この射は相対自己同型である。

 \(\mathcal L_d=\bigotimes_{j=0}^{d-1}(\sigma^j)^*\OO_E(1)\) は、
 滑らかなファイバーでは次数 \(3d\)、三角形では各成分次数 \(d\) を持つ。
 \(d>0\) なら \(H^1(E_h,\mathcal L_d)=0\)、\(h^0=3d\) であり、
 相対切断加群はlocally freeで底変更と両立する。
 自然な次数付き写像 \(S_h\to\bigoplus_dH^0(E_h,\mathcal L_d)\) の下で、
 \(G_h\) は一般の滑らかなファイバーで零なので、局所自由性から恒等的に零となる。
 三角形のtwisted homogeneous coordinate ringが次数一で生成されることを確認する。
 \(L\) のbasepoint-free pencilに対するpencil trickを使うと、
 多次数 \((d,d,d)\) の \(M\) について
 \(H^0(L)\otimes H^0(M)\to H^0(LM)\) は
 \(H^1(ML^{-1})=0\) なら全射である。
 \(d\ge2\) では各成分上の正値性からこの消滅が従う。
 \(d=1\) で \(ML^{-1}\) が非自明なら、多次数零の直線束の切断は
 三角形の三成分の接合条件により零で、Euler数も零なので同じ消滅が成り立つ。
 \(M\simeq L\) なら、平面三次曲線の斉次イデアルに二次式がないことから全射である。
 従って \(S_{h_0}/(G_{h_0})\to B(E_{h_0},L,\sigma_{h_0})\) は全射である。
 補題\ref{lem:center}の正則性から左辺のHilbert係数は \(1,3d\) であり、
 右辺の係数と一致する。よってこの全射は同型で、\eqref{eq:triangle}が従う。
 \(\mathcal L_d\) の各成分上の次数が \(d\) とともに増大するので
 \(\sigma\)-amplenessも成り立ち、tailsは \(\operatorname{coh}E_{h_0}\) と同値になる。

 Gorenstein parameter零のOrlov同値と \cite[Theorem 5.8]{CCKM}から
 \[
 \MF^\ZZ(S_{h_0},G_{h_0})\simeq D^b(\operatorname{coh}E_{h_0}).
 \]
 三角形の節点 \(p\) のskyscraper sheafは、任意に高い次数で非零の自己Extを持つ。
 節点の完備局所環 \(\CC[[s,t]]/(st)\) の剰余体の最小自由分解は無限であり、
 局所大域Extスペクトル系列では高次Extの支持が \(p\) だけなので、
 この非零性は大域自己Extにも残る。従ってこのdg圏はproperでない。

 \(m=0\) の場合の非propernessはこれで示された。以下では \(m\ge1\) とする。
 \(P=\sum_i u_i^3\) とし、第\ref{sec:TS}節の二ブロック次数 \(\Gamma\) と
 Koszul factorization \(K_P\) を使う。
 \(K_P\) の全微分行列は \((u_1,\ldots,u_m)\) に属する。
 もし恒等写像が零ホモトピーなら、係数を原点で評価して非零自由加群上の恒等写像が
 零となるので矛盾する。従って \(H^0\REnd(K_P)\ne0\) である。
 節点の対象を \(N\) とすると、外部積の完全忠実性と体上のKünneth公式により
 \[
 H^r\REnd(N)\otimes H^0\REnd(K_P)
 \subset H^r\REnd(N\boxtimes K_P)
 \]
 が直和因子として残る。よって \(\Gamma\)-次数の全体圏はnonproperである。
 もし通常の整数次数付き圏がproperなら、その有限群 \(H_\Gamma\simeq\mu_3\) による
 同変圏も、Hom複体の不変部分が直和因子となることからproperとなる。
 これは式\eqref{eq:eq}に反する。

 さらに、この議論は単なるHom総次元の非有限性より強い。
 \(\Gamma\)-次数の一つの対象に、任意に高い正次数の自己Homが残っている。
 有限群同変化から忘却すると、同変Homは通常のHomの直和因子なので、
 通常次数でも同じ正次数非消滅を持つ一つの対象が得られる。
 \(m>0\) ではOrlovの完全忠実windowによりこれを \(D^b(\qgr B_{h_0,m})\) へ送れ、
 \(m=0\) では上の三角形との同値を用いる。
 もし \(\qgr B_{h_0,m}\) のhomological dimensionが有限なら、
 その有界導来圏の任意の固定された二対象の間の正次数Homは十分高次で消える。
 これは有界複体のcohomologyとExtからなる有限スペクトル系列で分かる。
 今得た対象はこれに反するから、homological dimensionは無限である。

 一方、周囲の \(S_{h_0}[u]\) はAS-regularであり、三次の中心元の正則性は
 \(m>0\) なら最後の中心変数についてmonicであること、\(m=0\) なら
 補題\ref{lem:center}から従う。従ってHilbert seriesは依然として
 \((1-t^3)/(1-t)^{m+3}\) であり、degree shiftの三乗も \([2]\) である。
 滑らかな射影多様体のadmissible成分はproperなので、求める非存在が従う。
\end{proof}

この例は、量子化を構成する開集合をさらに縮めて例外値を除くことを妨げない。
従って、零を含む十分小さな開集合での存在問題に対する反例ではない。

\subsection{量子化原稿への適用}

三次の中心元によるAS-regularな周囲の商
\[
 B_\hbar=A_\hbar/(\Omega_\hbar),\qquad B_0=\CC[x]/(F)
\]
を、次数ごとにflatな族として構成できていれば、
その形式完備化には定理\ref{thm:main}が適用できる。
この適用にはPoisson構造の分類を再度用いず、各Artin次数でのgraded flatnessだけを使う。
従って、量子化の構成自体の証明と、本稿の相対残余圏・復元の証明は独立に検証できる。
論文Aの代数的量子化定理への適用では、そこで構成する次数ごとのflatnessを用いる。
代数の量子化と圏の復元は異なる段階であり、それぞれの仮定を区別する。

形式的な可換モデルは、半古典極限の固定された \(X\) と同じとは限らない。
定理\ref{thm:global}の代表式では \(h^2xyz\) が現れる。
このモデルを初期の \(F\) と取り違えない。

\subsection{一般の有限パラメータと次の核心}

定理\ref{thm:etale}により、smooth properな代数的pairの族が与えられれば、
可換cubic点のétale近傍で存在問題は解ける。
一般の量子化族に残る最初の課題は、その相対saturationの構成である。
各Artin切断のsmooth proper性だけから、有限表示の代数的模型の存在を結論しない。

独立な複数ブロックでは、命題\ref{prop:multiblock-embedding}で得た完全忠実な像が
全MF圏を生成するかが、二つ以上の無限位数因子の場合の核心である。
命題\ref{prop:block-support}と条件\eqref{eq:block-generation}は、その具体的な帰着を与える。
高々一つの無限位数因子を持つ場合は定理\ref{thm:multiblock-model}で解決した。
中心三次式内の混合項は定理\ref{thm:global}と定理\ref{thm:multiblock-model}に含まれるが、
非可換ブロックと中心変数の混合項は含まれない。
また、étale局所モデルの元の底への降下と、実現開部分の外のファイバーの分類は
別の大域問題として残る。

